
\Data\ID{1003035901}
\Updated{2020-07-15 03:07:12}
\Level{A}
\SubArea{Limit of a sequence}
\LANGen{
\Question{
Find the limit \( \lim\limits_{n\rightarrow\infty}\left(4+\frac3{2n-1}\right) \).}
{\( 4 \)}{\( \frac32 \)}{\( \infty \)}{\( 0 \)}{\( \frac{11}2 \)}{}}
\LANGpl{
\Question{
Wyznacz granicę \( \lim\limits_{n\rightarrow\infty}\left(4+\frac3{2n-1}\right) \).}
{\( 4 \)}{\( \frac32 \)}{\( \infty \)}{\( 0 \)}{\( \frac{11}2 \)}{}}
\LANGcs{
\Question{
\( \lim\limits_{n\rightarrow\infty}\left(4+\frac3{2n-1}\right) \) je rovna:}
{\( 4 \)}{\( \frac32 \)}{\( \infty \)}{\( 0 \)}{\( \frac{11}2 \)}{}}
\LANGsk{
\Question{
\( \lim\limits_{n\rightarrow\infty}\left(4+\frac3{2n-1}\right) \) je rovná:}
{\( 4 \)}{\( \frac32 \)}{\( \infty \)}{\( 0 \)}{\( \frac{11}2 \)}{}}
\LANGes{
\Question{
Halla \( \lim\limits_{n\rightarrow\infty}\left(4+\frac3{2n-1}\right) \).}
{\( 4 \)}{\( \frac32 \)}{\( \infty \)}{\( 0 \)}{\( \frac{11}2 \)}{}}
\End

\Data\ID{1003035903}
\Updated{2020-07-15 03:07:12}
\Level{A}
\SubArea{Limit of a sequence}
\LANGen{
\Question{
Find the limit \( \lim\limits_{n\rightarrow\infty}\left(5n^2+2n-1\right) \).}
{\( \infty \)}{\( 5 \)}{\( 0 \)}{\( 2 \)}{\( -1 \)}{}}
\LANGpl{
\Question{
Wyznacz granicę \( \lim\limits_{n\rightarrow\infty}\left(5n^2+2n-1\right) \).}
{\( \infty \)}{\( 5 \)}{\( 0 \)}{\( 2 \)}{\( -1 \)}{}}
\LANGcs{
\Question{
\( \lim\limits_{n\rightarrow\infty}\left(5n^2+2n-1\right) \) je rovna:}
{\( \infty \)}{\( 5 \)}{\( 0 \)}{\( 2 \)}{\( -1 \)}{}}
\LANGsk{
\Question{
\( \lim\limits_{n\rightarrow\infty}\left(5n^2+2n-1\right) \) je rovná:}
{\( \infty \)}{\( 5 \)}{\( 0 \)}{\( 2 \)}{\( -1 \)}{}}
\LANGes{
\Question{
Halla \( \lim\limits_{n\rightarrow\infty}\left(5n^2+2n-1\right) \).}
{\( \infty \)}{\( 5 \)}{\( 0 \)}{\( 2 \)}{\( -1 \)}{}}
\End

\Data\ID{1003035904}
\Updated{2020-07-15 03:07:12}
\Level{A}
\SubArea{Limit of a sequence}
\LANGen{
\Question{
Find the limit \( \lim\limits_{n\rightarrow\infty}\left(\frac{2n-2}{n+1}\cdot\frac{5n+3}{n-4}\right) \).}
{\( 10 \)}{\( \infty \)}{\( 0 \)}{\( 7 \)}{\( -6 \)}{}}
\LANGpl{
\Question{
Wyznacz granicę \( \lim\limits_{n\rightarrow\infty}\left(\frac{2n-2}{n+1}\cdot\frac{5n+3}{n-4}\right) \).}
{\( 10 \)}{\( \infty \)}{\( 0 \)}{\( 7 \)}{\( -6 \)}{}}
\LANGcs{
\Question{
\( \lim\limits_{n\rightarrow\infty}\left(\frac{2n-2}{n+1}\cdot\frac{5n+3}{n-4}\right) \) je rovna:}
{\( 10 \)}{\( \infty \)}{\( 0 \)}{\( 7 \)}{\( -6 \)}{}}
\LANGsk{
\Question{
\( \lim\limits_{n\rightarrow\infty}\left(\frac{2n-2}{n+1}\cdot\frac{5n+3}{n-4}\right) \) je rovná:}
{\( 10 \)}{\( \infty \)}{\( 0 \)}{\( 7 \)}{\( -6 \)}{}}
\LANGes{
\Question{
Halla \( \lim\limits_{n\rightarrow\infty}\left(\frac{2n-2}{n+1}\cdot\frac{5n+3}{n-4}\right) \).}
{\( 10 \)}{\( \infty \)}{\( 0 \)}{\( 7 \)}{\( -6 \)}{}}
\End

\Data\ID{1003035905}
\Updated{2020-07-15 03:07:12}
\Level{A}
\SubArea{Limit of a sequence}
\LANGen{
\Question{
Find the limit \( \lim\limits_{n\rightarrow\infty} \left( \frac{2n-2}{n+1}+\frac{5n+3}{n-4} \right) \).}
{\( 7 \)}{\( \infty \)}{\( 0 \)}{\( 10 \)}{\( -6 \)}{}}
\LANGpl{
\Question{
Wyznacz granicę \( \lim\limits_{n\rightarrow\infty} \left( \frac{2n-2}{n+1}+\frac{5n+3}{n-4} \right) \).}
{\( 7 \)}{\( \infty \)}{\( 0 \)}{\( 10 \)}{\( -6 \)}{}}
\LANGcs{
\Question{
\( \lim\limits_{n\rightarrow\infty} \left( \frac{2n-2}{n+1}+\frac{5n+3}{n-4} \right) \) je rovna:}
{\( 7 \)}{\( \infty \)}{\( 0 \)}{\( 10 \)}{\( -6 \)}{}}
\LANGsk{
\Question{
\( \lim\limits_{n\rightarrow\infty} \left( \frac{2n-2}{n+1}+\frac{5n+3}{n-4} \right) \) je rovná:}
{\( 7 \)}{\( \infty \)}{\( 0 \)}{\( 10 \)}{\( -6 \)}{}}
\LANGes{
\Question{
Halla \( \lim\limits_{n\rightarrow\infty} \left( \frac{2n-2}{n+1}+\frac{5n+3}{n-4} \right) \).}
{\( 7 \)}{\( \infty \)}{\( 0 \)}{\( 10 \)}{\( -6 \)}{}}
\End

\Data\ID{1003035906}
\Updated{2020-07-15 03:07:12}
\Level{A}
\SubArea{Limit of a sequence}
\LANGen{
\Question{
Find the limit \( \lim\limits_{n\rightarrow\infty}\left( \frac4{n^3} +2+\frac{3n^2+5}{1-n^2}\right) \).}
{\( -1 \)}{\( 0 \)}{\( \infty \)}{\( 9 \)}{\( -\frac13 \)}{}}
\LANGpl{
\Question{
Wyznacz granicę \( \lim\limits_{n\rightarrow\infty}\left( \frac4{n^3} +2+\frac{3n^2+5}{1-n^2}\right) \).}
{\( -1 \)}{\( 0 \)}{\( \infty \)}{\( 9 \)}{\( -\frac13 \)}{}}
\LANGcs{
\Question{
\( \lim\limits_{n\rightarrow\infty}\left( \frac4{n^3} +2+\frac{3n^2+5}{1-n^2}\right) \) je rovna:}
{\( -1 \)}{\( 0 \)}{\( \infty \)}{\( 9 \)}{\( -\frac13 \)}{}}
\LANGsk{
\Question{
\( \lim\limits_{n\rightarrow\infty}\left( \frac4{n^3} +2+\frac{3n^2+5}{1-n^2}\right) \) je rovná:}
{\( -1 \)}{\( 0 \)}{\( \infty \)}{\( 9 \)}{\( -\frac13 \)}{}}
\LANGes{
\Question{
Halla \( \lim\limits_{n\rightarrow\infty}\left( \frac4{n^3} +2+\frac{3n^2+5}{1-n^2}\right) \).}
{\( -1 \)}{\( 0 \)}{\( \infty \)}{\( 9 \)}{\( -\frac13 \)}{}}
\End

\Data\ID{1003047301}
\Updated{2020-07-15 03:07:19}
\Level{A}
\SubArea{Limit of a sequence}
\LANGen{
\Question{
Which of the following expressions describes the correct calculation of the limit?
\[ L=\lim\limits_{n\to\infty}\frac{n^3+2n-3}{2n^3+5}\]}
{\( L=\lim\limits_{n\to\infty}\frac{1+\frac2{n^2}-\frac3{n^3}}{2+\frac{5}{n^3}}= \frac12 \)}{\( L=\frac{\infty^3+2\cdot\infty-3}{2\cdot\infty^3+5}=\infty \)}{\( L=\frac{\infty^3+2\cdot\infty-3}{2\cdot\infty^3+5}=0 \)}{\( L=\lim\limits_{n\to\infty}\frac{n\left(n^2+2\right)-3}{2n^3+5}= -\frac35 \)}{\( L=\lim\limits_{n\to\infty}\frac{\left(n^2+3\right)(n-1)}{2\left(n^3+\frac52\right)}= 0 \)}{}}
\LANGpl{
\Question{
Które z poniższych wyrażeń opisuje poprawne obliczenie granicy?
\[ L=\lim\limits_{n\to\infty}\frac{n^3+2n-3}{2n^3+5} \]}
{\( L=\lim\limits_{n\to\infty}\frac{1+\frac2{n^2}-\frac3{n^3}}{2+\frac{5}{n^3}}= \frac12 \)}{\( L=\frac{\infty^3+2\cdot\infty-3}{2\cdot\infty^3+5}=\infty \)}{\( L=\frac{\infty^3+2\cdot\infty-3}{2\cdot\infty^3+5}=0 \)}{\( L=\lim\limits_{n\to\infty}\frac{n\left(n^2+2\right)-3}{2n^3+5}= -\frac35 \)}{\( L=\lim\limits_{n\to\infty}\frac{\left(n^2+3\right)(n-1)}{2\left(n^3+\frac52\right)}= 0 \)}{}}
\LANGcs{
\Question{
Vyberte správný výpočet limity posloupnosti.
\[ L=\lim\limits_{n\to\infty}\frac{n^3+2n-3}{2n^3+5} \]}
{\( L=\lim\limits_{n\to\infty}\frac{1+\frac2{n^2}-\frac3{n^3}}{2+\frac{5}{n^3}}= \frac12 \)}{\( L=\frac{\infty^3+2\cdot\infty-3}{2\cdot\infty^3+5}=\infty \)}{\( L=\frac{\infty^3+2\cdot\infty-3}{2\cdot\infty^3+5}=0 \)}{\( L=\lim\limits_{n\to\infty}\frac{n\left(n^2+2\right)-3}{2n^3+5}= -\frac35 \)}{\( L=\lim\limits_{n\to\infty}\frac{\left(n^2+3\right)(n-1)}{2\left(n^3+\frac52\right)}= 0 \)}{}}
\LANGsk{
\Question{
Vyberte správny výpočet limity postupnosti.
\[ L=\lim\limits_{n\to\infty}\frac{n^3+2n-3}{2n^3+5}\]}
{\( L=\lim\limits_{n\to\infty}\frac{1+\frac2{n^2}-\frac3{n^3}}{2+\frac{5}{n^3}}= \frac12 \)}{\( L=\frac{\infty^3+2\cdot\infty-3}{2\cdot\infty^3+5}=\infty \)}{\( L=\frac{\infty^3+2\cdot\infty-3}{2\cdot\infty^3+5}=0 \)}{\( L=\lim\limits_{n\to\infty}\frac{n\left(n^2+2\right)-3}{2n^3+5}= -\frac35 \)}{\( L=\lim\limits_{n\to\infty}\frac{\left(n^2+3\right)(n-1)}{2\left(n^3+\frac52\right)}= 0 \)}{}}
\LANGes{
\Question{
Elige la expresión correcta para calcular el límite de la siguiente sucesión:
\[ L=\lim\limits_{n\to\infty}\frac{n^3+2n-3}{2n^3+5}\]}
{\( L=\lim\limits_{n\to\infty}\frac{1+\frac2{n^2}-\frac3{n^3}}{2+\frac{5}{n^3}}= \frac12 \)}{\( L=\frac{\infty^3+2\cdot\infty-3}{2\cdot\infty^3+5}=\infty \)}{\( L=\frac{\infty^3+2\cdot\infty-3}{2\cdot\infty^3+5}=0 \)}{\( L=\lim\limits_{n\to\infty}\frac{n\left(n^2+2\right)-3}{2n^3+5}= -\frac35 \)}{\( L=\lim\limits_{n\to\infty}\frac{\left(n^2+3\right)(n-1)}{2\left(n^3+\frac52\right)}= 0 \)}{}}
\End

\Data\ID{1003047302}
\Updated{2021-05-03 08:05:58}
\Level{A}
\SubArea{Limit of a sequence}
\LANGen{
\Question{
Choose the best first step to take in order to evaluate the limit of the following sequence.
\[ \left( \frac{4n^5+n^4-n^3+2}{7n^4-2n^2+7n} \right)^{\infty}_{n=1} \]}
{We factor out \( n^4 \) in the numerator and in the denominator separately.}{We factor out \( n^5 \) in the numerator and in the denominator separately.}{We factorize a polynomial in the denominator.}{We divide the denominator by \( n^4 \).}{We divide the numerator by \( n^5 \).}{}}
\LANGpl{
\Question{
Który z poniższych kroków powinien być pierwszy, by obliczyć granicę ciągu?
\[ \left( \frac{4n^5+n^4-n^3+2}{7n^4-2n^2+7n} \right)^{\infty}_{n=1} \]}
{Usuwamy \( n^4 \) osobno z licznika i z mianownika.}{Usuwamy  \( n^5 \) osobno z licznika i mianownika.}{Rozkładamy wielomian w mianowniku na czynniki.}{Dzielimy mianownik przez \( n^4 \).}{Dzielimy licznik przez \( n^5 \).}{}}
\LANGcs{
\Question{
Vyberte nejvhodnější první krok k úpravě a výpočtu limity posloupnosti.
\[ \left( \frac{4n^5+n^4-n^3+2}{7n^4-2n^2+7n} \right)^{\infty}_{n=1}\ \]}
{Vytkneme v čitateli i jmenovateli \( n^4 \).}{Vytkneme v čitateli i jmenovateli \( n^5 \).}{Rozložíme jmenovatel na součin.}{Vydělíme jmenovatel \( n^4 \).}{Vydělíme čitatel \( n^5 \).}{}}
\LANGsk{
\Question{
Vyberte najvhodnejší prvý krok k úprave a výpočtu limity postupnosti.
\[ \left( \frac{4n^5+n^4-n^3+2}{7n^4-2n^2+7n} \right)^{\infty}_{n=1}\ \]}
{Vyberieme v čitateli a menovateli \( n^4 \).}{Vyberieme v čitateli a menovateli \( n^5 \).}{Rozložíme menovateľ na súčin.}{Vydelíme menovateľ \( n^4 \).}{Vydelíme čitateľ \( n^5 \).}{}}
\LANGes{
\Question{
Elige cual sería el primer paso más adecuado para calcular el límite de la siguiente sucesión:
\[ \left( \frac{4n^5+n^4-n^3+2}{7n^4-2n^2+7n} \right)^{\infty}_{n=1} \]}
{Factorizar \( n^4 \) en el numerador y en el denominador.}{Factorizar \( n^5 \) en el numerador y en el denominador.}{Factorizar un polinomio en el denominador.}{Dividir el denominador por \( n^4 \).}{Dividir el numerador por \( n^5 \).}{}}
\End

\Data\ID{1003047303}
\Updated{2020-07-15 03:07:19}
\Level{A}
\SubArea{Limit of a sequence}
\LANGen{
\Question{
Find the limit.
\[ \lim_{n\to\infty}\frac{-3n^2+n-1}{9n^5-3n^2+3} \]}
{\( 0 \)}{\( -\frac13 \)}{\( \infty \)}{\( -\infty \)}{\( \frac13 \)}{}}
\LANGpl{
\Question{
Wyznacz granicę. \[ \lim_{n\to\infty}\frac{-3n^2+n-1}{9n^5-3n^2+3} \]}
{\( 0 \)}{\( -\frac13 \)}{\( \infty \)}{\( -\infty \)}{\( \frac13 \)}{}}
\LANGcs{
\Question{
Vypočítejte limitu.
\[ \lim_{n\to\infty}\frac{-3n^2+n-1}{9n^5-3n^2+3} \]}
{\( 0 \)}{\( -\frac13 \)}{\( \infty \)}{\( -\infty \)}{\( \frac13 \)}{}}
\LANGsk{
\Question{
Vypočítajte limitu.
\[ \lim_{n\to\infty}\frac{-3n^2+n-1}{9n^5-3n^2+3} \]}
{\( 0 \)}{\( -\frac13 \)}{\( \infty \)}{\( -\infty \)}{\( \frac13 \)}{}}
\LANGes{
\Question{
Halla
\[ \lim_{n\to\infty}\frac{-3n^2+n-1}{9n^5-3n^2+3} \]}
{\( 0 \)}{\( -\frac13 \)}{\( \infty \)}{\( -\infty \)}{\( \frac13 \)}{}}
\End

\Data\ID{1003047304}
\Updated{2020-07-15 03:07:19}
\Level{A}
\SubArea{Limit of a sequence}
\LANGen{
\Question{
Find the limit.
\[ \lim\limits_{n\to\infty}\frac{-n^2-3}{7n} \]}
{\( -\infty \)}{\( 0 \)}{\( \infty \)}{\( -\frac37 \)}{\( 1 \)}{}}
\LANGpl{
\Question{
Wyznacz granicę.
\[ \lim\limits_{n\to\infty}\frac{-n^2-3}{7n} \]}
{\( -\infty \)}{\( 0 \)}{\( \infty \)}{\( -\frac37 \)}{\( 1 \)}{}}
\LANGcs{
\Question{
Vypočítejte limitu.
\[ \lim\limits_{n\to\infty}\frac{-n^2-3}{7n} \]}
{\( -\infty \)}{\( 0 \)}{\( \infty \)}{\( -\frac37 \)}{\( 1 \)}{}}
\LANGsk{
\Question{
Vypočítajte limitu.
\[ \lim\limits_{n\to\infty}\frac{-n^2-3}{7n} \]}
{\( -\infty \)}{\( 0 \)}{\( \infty \)}{\( -\frac37 \)}{\( 1 \)}{}}
\LANGes{
\Question{
Halla:
\[ \lim\limits_{n\to\infty}\frac{-n^2-3}{7n} \]}
{\( -\infty \)}{\( 0 \)}{\( \infty \)}{\( -\frac37 \)}{\( 1 \)}{}}
\End

\Data\ID{1003047305}
\Updated{2020-07-15 03:07:19}
\Level{A}
\SubArea{Limit of a sequence}
\LANGen{
\Question{
The sequence 
\[ \left(\frac{12n^3+5n+1}{2n^3-6}\right)_{n=1}^{\infty} \]}
{is convergent and \( \lim\limits_{n\to\infty}\frac{12n^3+5n+1}{2n^3-6}=6 \).}{is convergent and \( \lim\limits_{n\to\infty}\frac{12n^3+5n+1}{2n^3-6}=0 \).}{is convergent and \( \lim\limits_{n\to\infty}\frac{12n^3+5n+1}{2n^3-6}=12 \).}{is divergent and \( \lim\limits_{n\to\infty}\frac{12n^3+5n+1}{2n^3-6}=\infty \).}{has no limit.}{}}
\LANGpl{
\Question{
Ciąg \[ \left(\frac{12n^3+5n+1}{2n^3-6}\right)_{n=1}^{\infty} \]}
{jest zbieżny i  \( \lim\limits_{n\to\infty}\frac{12n^3+5n+1}{2n^3-6}=6 \).}{jest zbieżny i  \( \lim\limits_{n\to\infty}\frac{12n^3+5n+1}{2n^3-6}=0 \).}{jest zbieżny i  \( \lim\limits_{n\to\infty}\frac{12n^3+5n+1}{2n^3-6}=12 \).}{jest rozbieżny i \( \lim\limits_{n\to\infty}\frac{12n^3+5n+1}{2n^3-6}=\infty \).}{nie ma granicy.}{}}
\LANGcs{
\Question{
Posloupnost
\[ \left(\frac{12n^3+5n+1}{2n^3-6}\right)_{n=1}^{\infty} \]}
{je konvergentní a platí: \( \lim\limits_{n\to\infty}\frac{12n^3+5n+1}{2n^3-6}=6 \).}{je konvergentní a platí: \( \lim\limits_{n\to\infty}\frac{12n^3+5n+1}{2n^3-6}=0 \).}{je konvergentní a platí: \( \lim\limits_{n\to\infty}\frac{12n^3+5n+1}{2n^3-6}=12 \).}{je divergentní a platí: \( \lim\limits_{n\to\infty}\frac{12n^3+5n+1}{2n^3-6}=\infty \).}{nemá limitu.}{}}
\LANGsk{
\Question{
Postupnosť
\[ \left(\frac{12n^3+5n+1}{2n^3-6}\right)_{n=1}^{\infty} \]}
{je konvergentná a platí: \( \lim\limits_{n\to\infty}\frac{12n^3+5n+1}{2n^3-6}=6 \).}{je konvergentná a platí: \( \lim\limits_{n\to\infty}\frac{12n^3+5n+1}{2n^3-6}=0 \).}{je konvergentná a platí: \( \lim\limits_{n\to\infty}\frac{12n^3+5n+1}{2n^3-6}=12 \).}{je divergentná a platí: \( \lim\limits_{n\to\infty}\frac{12n^3+5n+1}{2n^3-6}=\infty \).}{nemá limitu.}{}}
\LANGes{
\Question{
La sucesión
\[ \left(\frac{12n^3+5n+1}{2n^3-6}\right)_{n=1}^{\infty} \]}
{es convergente y \( \lim\limits_{n\to\infty}\frac{12n^3+5n+1}{2n^3-6}=6 \).}{es convergente y \( \lim\limits_{n\to\infty}\frac{12n^3+5n+1}{2n^3-6}=0 \).}{es convergente y \( \lim\limits_{n\to\infty}\frac{12n^3+5n+1}{2n^3-6}=12 \).}{es divergente y \( \lim\limits_{n\to\infty}\frac{12n^3+5n+1}{2n^3-6}=\infty \).}{no tiene límite.}{}}
\End

\Data\ID{1003047306}
\Updated{2020-07-15 03:07:19}
\Level{A}
\SubArea{Limit of a sequence}
\LANGen{
\Question{
Which of the following expressions describes the correct calculation of the limit?
\[ L=\lim\limits_{n\to\infty}\frac{7n^4+6n^3-5n^2}{8n^5-7n^4+6} \]}
{\( L=\lim\limits_{n\to\infty}\frac{\frac7n+\frac6{n^2}-\frac5{n^3}}{8-\frac7n+\frac6{n^5}}=0 \)}{\( L=\lim\limits_{n\to\infty}\frac{7+\frac6n-\frac5{n^2}}{8n-7+\frac6{n^4} }=-1 \)}{\( L=\lim\limits_{n\to\infty}\frac{7+6-5n^2}{8n-7n+6}=-\infty \)}{\( L=\lim\limits_{n\to\infty}\frac{7+\frac6n-\frac5{n^2}}{8-\frac7n+\frac6{n^5}}=\frac78 \)}{\( L=\lim\limits_{n\to\infty}\frac{\frac7n+\frac6{n^2}-\frac5{n^3}}{8n-7+\frac6{n^4}}=0 \)}{}}
\LANGpl{
\Question{
Które z poniższych wyrażeń przedstawia poprawne obliczenia granicy ciągu?
\[ L=\lim\limits_{n\to\infty}\frac{7n^4+6n^3-5n^2}{8n^5-7n^4+6} \]}
{\( L=\lim\limits_{n\to\infty}\frac{\frac7n+\frac6{n^2}-\frac5{n^3}}{8-\frac7n+\frac6{n^5}}=0 \)}{\( L=\lim\limits_{n\to\infty}\frac{7+\frac6n-\frac5{n^2}}{8n-7+\frac6{n^4} }=-1 \)}{\( L=\lim\limits_{n\to\infty}\frac{7+6-5n^2}{8n-7n+6}=-\infty \)}{\( L=\lim\limits_{n\to\infty}\frac{7+\frac6n-\frac5{n^2}}{8-\frac7n+\frac6{n^5}}=\frac78 \)}{\( L=\lim\limits_{n\to\infty}\frac{\frac7n+\frac6{n^2}-\frac5{n^3}}{8n-7+\frac6{n^4}}=0 \)}{}}
\LANGcs{
\Question{
Vyberte správný výpočet limity posloupnosti.
\[ L=\lim\limits_{n\to\infty}\frac{7n^4+6n^3-5n^2}{8n^5-7n^4+6} \]}
{\( L=\lim\limits_{n\to\infty}\frac{\frac7n+\frac6{n^2}-\frac5{n^3}}{8-\frac7n+\frac6{n^5}}=0 \)}{\( L=\lim\limits_{n\to\infty}\frac{7+\frac6n-\frac5{n^2}}{8n-7+\frac6{n^4} }=-1 \)}{\( L=\lim\limits_{n\to\infty}\frac{7+6-5n^2}{8n-7n+6}=-\infty \)}{\( L=\lim\limits_{n\to\infty}\frac{7+\frac6n-\frac5{n^2}}{8-\frac7n+\frac6{n^5}}=\frac78 \)}{\( L=\lim\limits_{n\to\infty}\frac{\frac7n+\frac6{n^2}-\frac5{n^3}}{8n-7+\frac6{n^4}}=0 \)}{}}
\LANGsk{
\Question{
Vyberte správny výpočet limity postupnosti.
\[ L=\lim\limits_{n\to\infty}\frac{7n^4+6n^3-5n^2}{8n^5-7n^4+6} \]}
{\( L=\lim\limits_{n\to\infty}\frac{\frac7n+\frac6{n^2}-\frac5{n^3}}{8-\frac7n+\frac6{n^5}}=0 \)}{\( L=\lim\limits_{n\to\infty}\frac{7+\frac6n-\frac5{n^2}}{8n-7+\frac6{n^4} }=-1 \)}{\( L=\lim\limits_{n\to\infty}\frac{7+6-5n^2}{8n-7n+6}=-\infty \)}{\( L=\lim\limits_{n\to\infty}\frac{7+\frac6n-\frac5{n^2}}{8-\frac7n+\frac6{n^5}}=\frac78 \)}{\( L=\lim\limits_{n\to\infty}\frac{\frac7n+\frac6{n^2}-\frac5{n^3}}{8n-7+\frac6{n^4}}=0 \)}{}}
\LANGes{
\Question{
Elige la expresión correcta para calcular el límite de la siguiente sucesión:
\[ L=\lim\limits_{n\to\infty}\frac{7n^4+6n^3-5n^2}{8n^5-7n^4+6} \]}
{\( L=\lim\limits_{n\to\infty}\frac{\frac7n+\frac6{n^2}-\frac5{n^3}}{8-\frac7n+\frac6{n^5}}=0 \)}{\( L=\lim\limits_{n\to\infty}\frac{7+\frac6n-\frac5{n^2}}{8n-7+\frac6{n^4} }=-1 \)}{\( L=\lim\limits_{n\to\infty}\frac{7+6-5n^2}{8n-7n+6}=-\infty \)}{\( L=\lim\limits_{n\to\infty}\frac{7+\frac6n-\frac5{n^2}}{8-\frac7n+\frac6{n^5}}=\frac78 \)}{\( L=\lim\limits_{n\to\infty}\frac{\frac7n+\frac6{n^2}-\frac5{n^3}}{8n-7+\frac6{n^4}}=0 \)}{}}
\End

\Data\ID{1003047307}
\Updated{2020-07-15 03:07:19}
\Level{A}
\SubArea{Limit of a sequence}
\LANGen{
\Question{
Find the limit.
\[ \lim\limits_{n\to\infty}\frac{4n^2+n-2}{-5n^2-3n} \]}
{\( -\frac45 \)}{\( 0 \)}{\( \infty \)}{\( -\infty \)}{\( \frac45 \)}{}}
\LANGpl{
\Question{
Wyznacz granicę. \[ \lim\limits_{n\to\infty}\frac{4n^2+n-2}{-5n^2-3n} \]}
{\( -\frac45 \)}{\( 0 \)}{\( \infty \)}{\( -\infty \)}{\( \frac45 \)}{}}
\LANGcs{
\Question{
Vypočítejte limitu.
\[ \lim\limits_{n\to\infty}\frac{4n^2+n-2}{-5n^2-3n} \]}
{\( -\frac45 \)}{\( 0 \)}{\( \infty \)}{\( -\infty \)}{\( \frac45 \)}{}}
\LANGsk{
\Question{
Vypočítajte limitu.
\[ \lim\limits_{n\to\infty}\frac{4n^2+n-2}{-5n^2-3n} \]}
{\( -\frac45 \)}{\( 0 \)}{\( \infty \)}{\( -\infty \)}{\( \frac45 \)}{}}
\LANGes{
\Question{
Halla:
\[ \lim\limits_{n\to\infty}\frac{4n^2+n-2}{-5n^2-3n} \]}
{\( -\frac45 \)}{\( 0 \)}{\( \infty \)}{\( -\infty \)}{\( \frac45 \)}{}}
\End

\Data\ID{1003047308}
\Updated{2021-05-03 08:05:30}
\Level{A}
\SubArea{Limit of a sequence}
\LANGen{
\Question{
Choose the best first step to take in order to evaluate the limit of the following sequence.
\[ \left(\frac{3n^2-2n+4}{8n^2+13n+2}\right)_{n=1}^{\infty} \]}
{We divide the numerator and the denominator by \( n^2 \).}{We divide the numerator and the denominator by \( n \).}{We substitute \( n=\infty \).}{We factor out \( n \) in the numerator and in the denominator separately.}{We factor out \( 8 \) in the numerator and in the denominator separately.}{}}
\LANGpl{
\Question{
Które z poniższych działań jest najlepszym pierwszym krokiem do obliczenia granicy ciągu?
\[ \left(\frac{3n^2-2n+4}{8n^2+13n+2}\right)_{n=1}^{\infty}  \]}
{Dzielimy licznik i mianownik przez \( n^2 \).}{Dzielimy licznik i mianownik przez \( n \).}{Podstawiamy  \( n=\infty \).}{Usuwamy  \( n \) osobno z licznika i mianownika.}{Usuwamy \( 8 \) osobno z licznika i mianownika.}{}}
\LANGcs{
\Question{
Vyberte nejvhodnější první krok k úpravě a výpočtu limity posloupnosti.
\[ \left(\frac{3n^2-2n+4}{8n^2+13n+2}\right)_{n=1}^{\infty} \ \]}
{Vydělíme čitatel i jmenovatel \( n^2 \).}{Vydělíme čitatel i jmenovatel \( n \).}{Dosadíme \( n=\infty \).}{Vytkneme v čitateli i jmenovateli \( n \).}{Vytkneme v čitateli i jmenovateli \( 8 \).}{}}
\LANGsk{
\Question{
Vyberte najvhodnejší prvý krok k úprave a výpočtu limity postupnosti.
\[ \left(\frac{3n^2-2n+4}{8n^2+13n+2}\right)_{n=1}^{\infty} \ \]}
{Vydelíme čitateľ a menovateľ \( n^2 \).}{Vydelíme čitateľ a menovateľ \( n \).}{Dosadíme \( n=\infty \).}{Vyberieme v čitateli a menovateli \( n \).}{Vyberieme v čitateli i menovateli \( 8 \).}{}}
\LANGes{
\Question{
Elige el primer paso más adecuado para calcular el límite de la siguiente sucesión:
\[ \left(\frac{3n^2-2n+4}{8n^2+13n+2}\right)_{n=1}^{\infty} \]}
{Dividir el numerador y el denominador por \( n^2 \).}{Dividir el numerador y el denominador por \( n \).}{Sustituir \( n=\infty \).}{Factorizar \( n \) en el numerador y en el denominador.}{Factorizar \( 8 \) en el numerador y en el denominador.}{}}
\End

\Data\ID{1003047309}
\Updated{2020-07-15 03:07:19}
\Level{A}
\SubArea{Limit of a sequence}
\LANGen{
\Question{
The sequence
\[ \left(\frac{3n^5+2n^3+1}{n^3+3}\right)_{n=1}^{\infty} \]}
{is divergent and \( \lim\limits_{n\to\infty}\frac{3n^5+2n^3+1}{n^3+3}=\infty \).}{is convergent and \( \lim\limits_{n\to\infty}\frac{3n^5+2n^3+1}{n^3+3}=0 \).}{is convergent and \( \lim\limits_{n\to\infty}\frac{3n^5+2n^3+1}{n^3+3}=3 \).}{is divergent and \( \lim\limits_{n\to\infty}\frac{3n^5+2n^3+1}{n^3+3}=-\infty \).}{has no limit.}{}}
\LANGpl{
\Question{
Ciąg 
\[ \left(\frac{3n^5+2n^3+1}{n^3+3}\right)_{n=1}^{\infty} \]}
{jest rozbieżny i \( \lim\limits_{n\to\infty}\frac{3n^5+2n^3+1}{n^3+3}=\infty \).}{jest zbieżny i  \( \lim\limits_{n\to\infty}\frac{3n^5+2n^3+1}{n^3+3}=0 \).}{jest zbieżny i  \( \lim\limits_{n\to\infty}\frac{3n^5+2n^3+1}{n^3+3}=3 \).}{jest rozbieżny i  \( \lim\limits_{n\to\infty}\frac{3n^5+2n^3+1}{n^3+3}=-\infty \).}{nie ma granicy.}{}}
\LANGcs{
\Question{
Posloupnost
\[ \left(\frac{3n^5+2n^3+1}{n^3+3}\right)_{n=1}^{\infty} \]}
{je divergentní a platí: \( \lim\limits_{n\to\infty}\frac{3n^5+2n^3+1}{n^3+3}=\infty \).}{je konvergentní a platí: \( \lim\limits_{n\to\infty}\frac{3n^5+2n^3+1}{n^3+3}=0 \).}{je konvergentní a platí: \( \lim\limits_{n\to\infty}\frac{3n^5+2n^3+1}{n^3+3}=3 \).}{je divergentní a platí: \( \lim\limits_{n\to\infty}\frac{3n^5+2n^3+1}{n^3+3}=-\infty \).}{nemá limitu.}{}}
\LANGsk{
\Question{
Postupnosť
\[ \left(\frac{3n^5+2n^3+1}{n^3+3}\right)_{n=1}^{\infty} \]}
{je divergentná a platí: \( \lim\limits_{n\to\infty}\frac{3n^5+2n^3+1}{n^3+3}=\infty \).}{je konvergentná a platí: \( \lim\limits_{n\to\infty}\frac{3n^5+2n^3+1}{n^3+3}=0 \).}{je konvergentná a platí: \( \lim\limits_{n\to\infty}\frac{3n^5+2n^3+1}{n^3+3}=3 \).}{je divergentná a platí: \( \lim\limits_{n\to\infty}\frac{3n^5+2n^3+1}{n^3+3}=-\infty \).}{nemá limitu.}{}}
\LANGes{
\Question{
La sucesión
\[ \left(\frac{3n^5+2n^3+1}{n^3+3}\right)_{n=1}^{\infty} \]}
{es divergente y \( \lim\limits_{n\to\infty}\frac{3n^5+2n^3+1}{n^3+3}=\infty \).}{es convergente y \( \lim\limits_{n\to\infty}\frac{3n^5+2n^3+1}{n^3+3}=0 \).}{es convergente y \( \lim\limits_{n\to\infty}\frac{3n^5+2n^3+1}{n^3+3}=3 \).}{es divergente y \( \lim\limits_{n\to\infty}\frac{3n^5+2n^3+1}{n^3+3}=-\infty \).}{no tiene límite.}{}}
\End

\Data\ID{1003047310}
\Updated{2020-07-15 03:07:19}
\Level{A}
\SubArea{Limit of a sequence}
\LANGen{
\Question{
Find the limit.
\[ \lim\limits_{n\to\infty}\frac{-4n+5}{2n+2} \]}
{\( -2 \)}{\( 0 \)}{\( \infty \)}{\( -\infty \)}{\( \frac52 \)}{}}
\LANGpl{
\Question{
Wyznacz granicę.
\[ \lim\limits_{n\to\infty}\frac{-4n+5}{2n+2} \]}
{\( -2 \)}{\( 0 \)}{\( \infty \)}{\( -\infty \)}{\( \frac52 \)}{}}
\LANGcs{
\Question{
Vypočítejte limitu.
\[ \lim\limits_{n\to\infty}\frac{-4n+5}{2n+2} \]}
{\( -2 \)}{\( 0 \)}{\( \infty \)}{\( -\infty \)}{\( \frac52 \)}{}}
\LANGsk{
\Question{
Vypočítajte limitu.
\[ \lim\limits_{n\to\infty}\frac{-4n+5}{2n+2} \]}
{\( -2 \)}{\( 0 \)}{\( \infty \)}{\( -\infty \)}{\( \frac52 \)}{}}
\LANGes{
\Question{
Halla:
\[ \lim\limits_{n\to\infty}\frac{-4n+5}{2n+2} \]}
{\( -2 \)}{\( 0 \)}{\( \infty \)}{\( -\infty \)}{\( \frac52 \)}{}}
\End

\Data\ID{1003047504}
\Updated{2020-07-15 03:07:12}
\Level{A}
\SubArea{Limit of a sequence}
\LANGen{
\Question{
Find the following limit. \[ \lim\limits_{n\rightarrow\infty}\frac{3\pi}{6n^2-1} \]}
{\( 0 \)}{\( \infty \)}{\( 3\pi \)}{\( \frac12 \)}{\( \frac16 \)}{}}
\LANGpl{
\Question{
Wyznacz granicę. \[ \lim\limits_{n\rightarrow\infty}\frac{3\pi}{6n^2-1} \]}
{\( 0 \)}{\( \infty \)}{\( 3\pi \)}{\( \frac12 \)}{\( \frac16 \)}{}}
\LANGcs{
\Question{
\( \lim\limits_{n\rightarrow\infty}\frac{3\pi}{6n^2-1} \) je rovna:}
{\( 0 \)}{\( \infty \)}{\( 3\pi \)}{\( \frac12 \)}{\( \frac16 \)}{}}
\LANGsk{
\Question{
\( \lim\limits_{n\rightarrow\infty}\frac{3\pi}{6n^2-1} \) je rovná:}
{\( 0 \)}{\( \infty \)}{\( 3\pi \)}{\( \frac12 \)}{\( \frac16 \)}{}}
\LANGes{
\Question{
Halla: \[ \lim\limits_{n\rightarrow\infty}\frac{3\pi}{6n^2-1} \]}
{\( 0 \)}{\( \infty \)}{\( 3\pi \)}{\( \frac12 \)}{\( \frac16 \)}{}}
\End

\Data\ID{9000063601}
\Updated{2020-07-15 03:07:37}
\Level{A}
\SubArea{Limit of a sequence}
\LANGen{
\Question{
Find the following limit.
\[ 
 \lim _{n\to \infty }\frac{2n + 3} 
{3n - 2}
\]}
{\(\frac{2}
{3}\)}{\(-\frac{3} 
{2}\)}{\(0\)}{\(1\)}{}{}}
\LANGpl{
\Question{
Wyznacz granicę ciągu.
\[ 
 \lim _{n\to \infty }\frac{2n + 3} 
{3n - 2}
\]}
{\(\frac{2}
{3}\)}{\(-\frac{3} 
{2}\)}{\(0\)}{\(1\)}{}{}}
\LANGcs{
\Question{
\(\lim _{n\to \infty }\frac{2n+3} 
{3n-2}\) je
rovna:}
{\(\frac{2}
{3}\)}{\(-\frac{3} 
{2}\)}{\(0\)}{\(1\)}{}{}}
\LANGsk{
\Question{
\(\lim _{n\to \infty }\frac{2n+3} 
{3n-2}\) je
rovná:}
{\(\frac{2}
{3}\)}{\(-\frac{3} 
{2}\)}{\(0\)}{\(1\)}{}{}}
\LANGes{
\Question{
Halla:
\[ 
 \lim _{n\to \infty }\frac{2n + 3} 
{3n - 2}
\]}
{\(\frac{2}
{3}\)}{\(-\frac{3} 
{2}\)}{\(0\)}{\(1\)}{}{}}
\End

\Data\ID{9000063603}
\Updated{2020-07-15 03:07:37}
\Level{A}
\SubArea{Limit of a sequence}
\LANGen{
\Question{
Find the following limit.
\[ 
 \lim _{n\to \infty }\frac{2n^{2} + 1} 
 {3n - 1}
\]}
{\(\infty \)}{\(\frac{3}
{2}\)}{\(0\)}{\(- 1\)}{}{}}
\LANGpl{
\Question{
Wyznacz granicę ciągu.
\[ 
 \lim _{n\to \infty }\frac{2n^{2} + 1} 
 {3n - 1}
\]}
{\(\infty \)}{\(\frac{3}
{2}\)}{\(0\)}{\(- 1\)}{}{}}
\LANGcs{
\Question{
\(\lim _{n\to \infty }\frac{2n^{2}+1} 
 {3n-1} \) je
rovna:}
{\(\infty \)}{\(\frac{3}
{2}\)}{\(0\)}{\(- 1\)}{}{}}
\LANGsk{
\Question{
\(\lim _{n\to \infty }\frac{2n^{2}+1} 
 {3n-1} \) je
rovná:}
{\(\infty \)}{\(\frac{3}
{2}\)}{\(0\)}{\(- 1\)}{}{}}
\LANGes{
\Question{
Halla:
\[ 
 \lim _{n\to \infty }\frac{2n^{2} + 1} 
 {3n - 1}
\]}
{\(\infty \)}{\(\frac{3}
{2}\)}{\(0\)}{\(- 1\)}{}{}}
\End

\Data\ID{9000063606}
\Updated{2020-07-15 03:07:37}
\Level{A}
\SubArea{Limit of a sequence}
\LANGen{
\Question{
Find the following limit.
\[ 
 \lim _{n\to \infty }\frac{3n^{2} - 2n + 1} 
 {2n^{3} - 4}
\]}
{\(0\)}{\(\frac{3}
{2}\)}{\(\frac{1}
{2}\)}{\(-\frac{1} 
{4}\)}{}{}}
\LANGpl{
\Question{
Wyznacz granicę ciągu.
\[ 
 \lim _{n\to \infty }\frac{3n^{2} - 2n + 1} 
 {2n^{3} - 4}
\]}
{\(0\)}{\(\frac{3}
{2}\)}{\(\frac{1}
{2}\)}{\(-\frac{1} 
{4}\)}{}{}}
\LANGcs{
\Question{
\(\lim _{n\to \infty }\frac{3n^{2}-2n+1} 
 {2n^{3}-4} \) je
rovna:}
{\(0\)}{\(\frac{3}
{2}\)}{\(\frac{1}
{2}\)}{\(-\frac{1} 
{4}\)}{}{}}
\LANGsk{
\Question{
\(\lim _{n\to \infty }\frac{3n^{2}-2n+1} 
 {2n^{3}-4} \) je
rovná:}
{\(0\)}{\(\frac{3}
{2}\)}{\(\frac{1}
{2}\)}{\(-\frac{1} 
{4}\)}{}{}}
\LANGes{
\Question{
Halla:
\[ 
 \lim _{n\to \infty }\frac{3n^{2} - 2n + 1} 
 {2n^{3} - 4}
\]}
{\(0\)}{\(\frac{3}
{2}\)}{\(\frac{1}
{2}\)}{\(-\frac{1} 
{4}\)}{}{}}
\End

\Data\ID{9000063609}
\Updated{2020-07-15 03:07:37}
\Level{A}
\SubArea{Limit of a sequence}
\LANGen{
\Question{
Find the following limit.
\[ 
 \lim _{n\to \infty }\left ( \frac{n}
{n - 1} + \frac{n + 2} 
{n + 1}\right )
\]}
{\(2\)}{\(- 1\)}{\(0\)}{\(1\)}{}{}}
\LANGpl{
\Question{
Wyznacz granicę ciągu.
\[ 
 \lim _{n\to \infty }\left ( \frac{n}
{n - 1} + \frac{n + 2} 
{n + 1}\right )
\]}
{\(2\)}{\(- 1\)}{\(0\)}{\(1\)}{}{}}
\LANGcs{
\Question{
\(\lim _{n\to \infty }\left ( \frac{n}
{n-1} + \frac{n+2} 
{n+1}\right )\) je
rovna:}
{\(2\)}{\(- 1\)}{\(0\)}{\(1\)}{}{}}
\LANGsk{
\Question{
\(\lim _{n\to \infty }\left ( \frac{n}
{n-1} + \frac{n+2} 
{n+1}\right )\) je
rovná:}
{\(2\)}{\(- 1\)}{\(0\)}{\(1\)}{}{}}
\LANGes{
\Question{
Halla:
\[ 
 \lim _{n\to \infty }\left ( \frac{n}
{n - 1} + \frac{n + 2} 
{n + 1}\right )
\]}
{\(2\)}{\(- 1\)}{\(0\)}{\(1\)}{}{}}
\End

\Data\ID{1003035902}
\Updated{2020-07-15 03:07:12}
\Level{B}
\SubArea{Limit of a sequence}
\LANGen{
\Question{
Find the limit \( \lim\limits_{n\rightarrow\infty}(-1)^n\cdot\frac{3n}{4n^2+7} \).}
{\( 0 \)}{\( \frac34 \)}{\( -\frac34 \)}{\( \frac37 \)}{\( -\infty \)}{}}
\LANGpl{
\Question{
Wyznacz granicę \( \lim\limits_{n\rightarrow\infty}(-1)^n\cdot\frac{3n}{4n^2+7} \).}
{\( 0 \)}{\( \frac34 \)}{\( -\frac34 \)}{\( \frac37 \)}{\( -\infty \)}{}}
\LANGcs{
\Question{
\( \lim\limits_{n\rightarrow\infty}(-1)^n\cdot\frac{3n}{4n^2+7} \) je rovna:}
{\( 0 \)}{\( \frac34 \)}{\( -\frac34 \)}{\( \frac37 \)}{\( -\infty \)}{}}
\LANGsk{
\Question{
\( \lim\limits_{n\rightarrow\infty}(-1)^n\cdot\frac{3n}{4n^2+7} \) je rovná:}
{\( 0 \)}{\( \frac34 \)}{\( -\frac34 \)}{\( \frac37 \)}{\( -\infty \)}{}}
\LANGes{
\Question{
Halla \( \lim\limits_{n\rightarrow\infty}(-1)^n\cdot\frac{3n}{4n^2+7} \).}
{\( 0 \)}{\( \frac34 \)}{\( -\frac34 \)}{\( \frac37 \)}{\( -\infty \)}{}}
\End

\Data\ID{1003035907}
\Updated{2020-07-15 03:07:12}
\Level{B}
\SubArea{Limit of a sequence}
\LANGen{
\Question{
Find the limit of the sequence \( \left(\left( \frac32 \right)^n \right)_{n=1}^{\infty} \).}
{\( \lim\limits_{n\rightarrow\infty}\left( \frac32 \right)^n =\infty \)}{\( \lim\limits_{n\rightarrow\infty}\left( \frac32 \right)^n =\frac32 \)}{\( \lim\limits_{n\rightarrow\infty}\left( \frac32 \right)^n =\frac{81}{16} \)}{\( \lim\limits_{n\rightarrow\infty}\left( \frac32 \right)^n = 0\)}{\( \lim\limits_{n\rightarrow\infty}\left( \frac32 \right)^n \) does not exist.}{}}
\LANGpl{
\Question{
Wyznacz granicę dla ciągu \( \left(\left( \frac32 \right)^n \right)_{n=1}^{\infty} \).}
{\( \lim\limits_{n\rightarrow\infty}\left( \frac32 \right)^n =\infty \)}{\( \lim\limits_{n\rightarrow\infty}\left( \frac32 \right)^n =\frac32 \)}{\( \lim\limits_{n\rightarrow\infty}\left( \frac32 \right)^n =\frac{81}{16} \)}{\( \lim\limits_{n\rightarrow\infty}\left( \frac32 \right)^n = 0\)}{\( \lim\limits_{n\rightarrow\infty}\left( \frac32 \right)^n \) nie istnieje.}{}}
\LANGcs{
\Question{
Určete limitu posloupnosti \( \left(\left( \frac32 \right)^n \right)_{n=1}^{\infty} \).}
{\( \lim\limits_{n\rightarrow\infty}\left( \frac32 \right)^n =\infty \)}{\( \lim\limits_{n\rightarrow\infty}\left( \frac32 \right)^n =\frac32 \)}{\( \lim\limits_{n\rightarrow\infty}\left( \frac32 \right)^n =\frac{81}{16} \)}{\( \lim\limits_{n\rightarrow\infty}\left( \frac32 \right)^n = 0 \)}{Daná posloupnost nemá limitu.}{}}
\LANGsk{
\Question{
Určte limitu postupnosti \( \left(\left( \frac32 \right)^n \right)_{n=1}^{\infty} \).}
{\( \lim\limits_{n\rightarrow\infty}\left( \frac32 \right)^n =\infty \)}{\( \lim\limits_{n\rightarrow\infty}\left( \frac32 \right)^n =\frac32 \)}{\( \lim\limits_{n\rightarrow\infty}\left( \frac32 \right)^n =\frac{81}{16} \)}{\( \lim\limits_{n\rightarrow\infty}\left( \frac32 \right)^n = 0 \)}{Daná postupnosť nemá limitu.}{}}
\LANGes{
\Question{
Calcula el límite de la sucesión \( \left(\left( \frac32 \right)^n \right)_{n=1}^{\infty} \).}
{\( \lim\limits_{n\rightarrow\infty}\left( \frac32 \right)^n =\infty \)}{\( \lim\limits_{n\rightarrow\infty}\left( \frac32 \right)^n =\frac32 \)}{\( \lim\limits_{n\rightarrow\infty}\left( \frac32 \right)^n =\frac{81}{16} \)}{\( \lim\limits_{n\rightarrow\infty}\left( \frac32 \right)^n = 0\)}{\( \lim\limits_{n\rightarrow\infty}\left( \frac32 \right)^n \) no existe.}{}}
\End

\Data\ID{1003035908}
\Updated{2020-07-15 03:07:12}
\Level{B}
\SubArea{Limit of a sequence}
\LANGen{
\Question{
Find the limit of the sequence \( \left(\left( \frac23 \right)^n\right)_{n=1}^{\infty} \).}
{\( \lim\limits_{n\rightarrow\infty}\left( \frac23 \right)^n =0 \)}{\( \lim\limits_{n\rightarrow\infty}\left( \frac23 \right)^n =-\infty \)}{\( \lim\limits_{n\rightarrow\infty}\left( \frac23 \right)^n =\frac{16}{81} \)}{\( \lim\limits_{n\rightarrow\infty}\left( \frac23 \right)^n =\frac23 \)}{\( \lim\limits_{n\rightarrow\infty}\left( \frac23 \right)^n \)  does not exist.}{}}
\LANGpl{
\Question{
Wyznacz granicę dla ciągu \( \left(\left( \frac23 \right)^n\right)_{n=1}^{\infty} \).}
{\( \lim\limits_{n\rightarrow\infty}\left( \frac23 \right)^n =0 \)}{\( \lim\limits_{n\rightarrow\infty}\left( \frac23 \right)^n =-\infty \)}{\( \lim\limits_{n\rightarrow\infty}\left( \frac23 \right)^n =\frac{16}{81} \)}{\( \lim\limits_{n\rightarrow\infty}\left( \frac23 \right)^n =\frac23 \)}{\( \lim\limits_{n\rightarrow\infty}\left( \frac23 \right)^n \)  nie istnieje.}{}}
\LANGcs{
\Question{
Určete limitu posloupnosti \( \left(\left( \frac23 \right)^n\right)_{n=1}^{\infty} \).}
{\( \lim\limits_{n\rightarrow\infty}\left( \frac23 \right)^n =0 \)}{\( \lim\limits_{n\rightarrow\infty}\left( \frac23 \right)^n =-\infty \)}{\( \lim\limits_{n\rightarrow\infty}\left( \frac23 \right)^n =\frac{16}{81} \)}{\( \lim\limits_{n\rightarrow\infty}\left( \frac23 \right)^n =\frac23 \)}{Daná posloupnost nemá limitu.}{}}
\LANGsk{
\Question{
Určte limitu postupnosti \( \left(\left( \frac23 \right)^n\right)_{n=1}^{\infty} \).}
{\( \lim\limits_{n\rightarrow\infty}\left( \frac23 \right)^n =0 \)}{\( \lim\limits_{n\rightarrow\infty}\left( \frac23 \right)^n =-\infty \)}{\( \lim\limits_{n\rightarrow\infty}\left( \frac23 \right)^n =\frac{16}{81} \)}{\( \lim\limits_{n\rightarrow\infty}\left( \frac23 \right)^n =\frac23 \)}{Daná postupnosť nemá limitu.}{}}
\LANGes{
\Question{
Calcula el límite de la sucesión \( \left(\left( \frac23 \right)^n\right)_{n=1}^{\infty} \).}
{\( \lim\limits_{n\rightarrow\infty}\left( \frac23 \right)^n =0 \)}{\( \lim\limits_{n\rightarrow\infty}\left( \frac23 \right)^n =-\infty \)}{\( \lim\limits_{n\rightarrow\infty}\left( \frac23 \right)^n =\frac{16}{81} \)}{\( \lim\limits_{n\rightarrow\infty}\left( \frac23 \right)^n =\frac23 \)}{\( \lim\limits_{n\rightarrow\infty}\left( \frac23 \right)^n \)  no existe.}{}}
\End

\Data\ID{1003035909}
\Updated{2020-07-15 03:07:12}
\Level{B}
\SubArea{Limit of a sequence}
\LANGen{
\Question{
Find the limit of the sequence \( \left(\left( -\frac32 \right)^n \right)_{n=1}^{\infty} \).}
{\( \lim\limits_{n\rightarrow\infty}\left( -\frac32 \right)^n \) does not exist.}{\( \lim\limits_{n\rightarrow\infty}\left( -\frac32 \right)^n = \infty \)}{\( \lim\limits_{n\rightarrow\infty}\left( -\frac32 \right)^n = 0 \)}{\( \lim\limits_{n\rightarrow\infty}\left( -\frac32 \right)^n = -\infty \)}{\( \lim\limits_{n\rightarrow\infty}\left( -\frac32 \right)^n = -\frac32\)}{}}
\LANGpl{
\Question{
Wyznacz granicę dla ciągu \( \left(\left( -\frac32 \right)^n \right)_{n=1}^{\infty} \).}
{\( \lim\limits_{n\rightarrow\infty}\left( -\frac32 \right)^n \) nie istnieje.}{\( \lim\limits_{n\rightarrow\infty}\left( -\frac32 \right)^n = \infty \)}{\( \lim\limits_{n\rightarrow\infty}\left( -\frac32 \right)^n = 0 \)}{\( \lim\limits_{n\rightarrow\infty}\left( -\frac32 \right)^n = -\infty \)}{\( \lim\limits_{n\rightarrow\infty}\left( -\frac32 \right)^n = -\frac32\)}{}}
\LANGcs{
\Question{
Určete limitu posloupnosti \( \left(\left( -\frac32 \right)^n \right)_{n=1}^{\infty} \).}
{Daná posloupnost nemá limitu.}{\( \lim\limits_{n\rightarrow\infty}\left( -\frac32 \right)^n = \infty \)}{\( \lim\limits_{n\rightarrow\infty}\left( -\frac32 \right)^n = 0 \)}{\( \lim\limits_{n\rightarrow\infty}\left( -\frac32 \right)^n = -\infty \)}{\( \lim\limits_{n\rightarrow\infty}\left( -\frac32 \right)^n = -\frac32\)}{}}
\LANGsk{
\Question{
Určte limitu postupnosti \( \left(\left( -\frac32 \right)^n \right)_{n=1}^{\infty} \).}
{Daná postupnosť nemá limitu.}{\( \lim\limits_{n\rightarrow\infty}\left( -\frac32 \right)^n = \infty \)}{\( \lim\limits_{n\rightarrow\infty}\left( -\frac32 \right)^n = 0 \)}{\( \lim\limits_{n\rightarrow\infty}\left( -\frac32 \right)^n = -\infty \)}{\( \lim\limits_{n\rightarrow\infty}\left( -\frac32 \right)^n = -\frac32\)}{}}
\LANGes{
\Question{
Calcula el límite de la sucesión \( \left(\left( -\frac32 \right)^n \right)_{n=1}^{\infty} \).}
{\( \lim\limits_{n\rightarrow\infty}\left( -\frac32 \right)^n \) no existe.}{\( \lim\limits_{n\rightarrow\infty}\left( -\frac32 \right)^n = \infty \)}{\( \lim\limits_{n\rightarrow\infty}\left( -\frac32 \right)^n = 0 \)}{\( \lim\limits_{n\rightarrow\infty}\left( -\frac32 \right)^n = -\infty \)}{\( \lim\limits_{n\rightarrow\infty}\left( -\frac32 \right)^n = -\frac32\)}{}}
\End

\Data\ID{1003035910}
\Updated{2020-07-15 03:07:12}
\Level{B}
\SubArea{Limit of a sequence}
\LANGen{
\Question{
Find the limit of the sequence \( \left( \left( -\frac23 \right)^n \right)_{n=1}^{\infty} \).}
{\( \lim\limits_{n\rightarrow\infty}\left( -\frac23 \right)^n=0 \)}{\( \lim\limits_{n\rightarrow\infty}\left( -\frac23 \right)^n=-\infty \)}{\( \lim\limits_{n\rightarrow\infty}\left( -\frac23 \right)^n=-\frac23 \)}{\( \lim\limits_{n\rightarrow\infty}\left( -\frac23 \right)^n=-\frac32 \)}{\( \lim\limits_{n\rightarrow\infty}\left( -\frac23 \right)^n \)  does not exist.}{}}
\LANGpl{
\Question{
Wyznacz granicę dla ciągu  \( \left( \left( -\frac23 \right)^n \right)_{n=1}^{\infty} \).}
{\( \lim\limits_{n\rightarrow\infty}\left( -\frac23 \right)^n=0 \)}{\( \lim\limits_{n\rightarrow\infty}\left( -\frac23 \right)^n=-\infty \)}{\( \lim\limits_{n\rightarrow\infty}\left( -\frac23 \right)^n=-\frac23 \)}{\( \lim\limits_{n\rightarrow\infty}\left( -\frac23 \right)^n=-\frac32 \)}{\( \lim\limits_{n\rightarrow\infty}\left( -\frac23 \right)^n \)   nie istnieje.}{}}
\LANGcs{
\Question{
Určete limitu posloupnosti \( \left( \left( -\frac23 \right)^n \right)_{n=1}^{\infty} \).}
{\( \lim\limits_{n\rightarrow\infty}\left( -\frac23 \right)^n=0 \)}{\( \lim\limits_{n\rightarrow\infty}\left( -\frac23 \right)^n=-\infty \)}{\( \lim\limits_{n\rightarrow\infty}\left( -\frac23 \right)^n=-\frac23 \)}{\( \lim\limits_{n\rightarrow\infty}\left( -\frac23 \right)^n=-\frac32 \)}{Daná posloupnost nemá limitu.}{}}
\LANGsk{
\Question{
Určte limitu postupnosti \( \left( \left( -\frac23 \right)^n \right)_{n=1}^{\infty} \).}
{\( \lim\limits_{n\rightarrow\infty}\left( -\frac23 \right)^n=0 \)}{\( \lim\limits_{n\rightarrow\infty}\left( -\frac23 \right)^n=-\infty \)}{\( \lim\limits_{n\rightarrow\infty}\left( -\frac23 \right)^n=-\frac23 \)}{\( \lim\limits_{n\rightarrow\infty}\left( -\frac23 \right)^n=-\frac32 \)}{Daná postupnosť nemá limitu.}{}}
\LANGes{
\Question{
Calcula el límite de la sucesión \( \left( \left( -\frac23 \right)^n \right)_{n=1}^{\infty} \).}
{\( \lim\limits_{n\rightarrow\infty}\left( -\frac23 \right)^n=0 \)}{\( \lim\limits_{n\rightarrow\infty}\left( -\frac23 \right)^n=-\infty \)}{\( \lim\limits_{n\rightarrow\infty}\left( -\frac23 \right)^n=-\frac23 \)}{\( \lim\limits_{n\rightarrow\infty}\left( -\frac23 \right)^n=-\frac32 \)}{\( \lim\limits_{n\rightarrow\infty}\left( -\frac23 \right)^n \)  no existe.}{}}
\End

\Data\ID{1003047401}
\Updated{2020-07-15 03:07:12}
\Level{B}
\SubArea{Limit of a sequence}
\LANGen{
\Question{
Choose the correct formula to calculate the limit. \[ L=\lim\limits_{n\rightarrow\infty}\frac{3\cdot5^n+2\cdot6^n}{2\cdot5^n+4\cdot7^n } \]}
{\( L=\lim\limits_{n\rightarrow\infty} \frac{3\cdot\left(\frac57\right)^n+2\cdot\left(\frac67\right)^n}{2\cdot\left(\frac57\right)^n+4} =0 \)}{\( L=\lim\limits_{n\rightarrow\infty}\frac{3\cdot\left(\frac56\right)^n+2}{2\cdot\left(\frac57\right)^n+4}=\frac12 \)}{\( L=\lim\limits_{n\rightarrow\infty}\frac{3+2\cdot\left(\frac65\right)^n}{2+4\cdot\left(\frac75\right)^n } =\frac32 \)}{\( L=\frac{3\cdot5^{\infty}+2\cdot6^{\infty}}{2\cdot5^{\infty}+4\cdot7^{\infty}}=\infty \)}{\( L=\lim\limits_{n\rightarrow\infty}\frac{3\cdot\left(\frac57\right)^n+2\cdot\left(\frac67\right)^n}{2\cdot\left(\frac57\right)^n+4\cdot\left(\frac77\right)^n}=\frac56 \)}{}}
\LANGpl{
\Question{
Wybierz odpowiedni wzór, za pomocą którego można obliczyć granicę ciągu. \[ L=\lim\limits_{n\rightarrow\infty}\frac{3\cdot5^n+2\cdot6^n}{2\cdot5^n+4\cdot7^n } \]}
{\( L=\lim\limits_{n\rightarrow\infty} \frac{3\cdot\left(\frac57\right)^n+2\cdot\left(\frac67\right)^n}{2\cdot\left(\frac57\right)^n+4} =0 \)}{\( L=\lim\limits_{n\rightarrow\infty}\frac{3\cdot\left(\frac56\right)^n+2}{2\cdot\left(\frac57\right)^n+4}=\frac12 \)}{\( L=\lim\limits_{n\rightarrow\infty}\frac{3+2\cdot\left(\frac65\right)^n}{2+4\cdot\left(\frac75\right)^n } =\frac32 \)}{\( L=\frac{3\cdot5^{\infty}+2\cdot6^{\infty}}{2\cdot5^{\infty}+4\cdot7^{\infty}}=\infty \)}{\( L=\lim\limits_{n\rightarrow\infty}\frac{3\cdot\left(\frac57\right)^n+2\cdot\left(\frac67\right)^n}{2\cdot\left(\frac57\right)^n+4\cdot\left(\frac77\right)^n}=\frac56 \)}{}}
\LANGcs{
\Question{
Vyberte správný výpočet limity posloupnosti. \[ L=\lim\limits_{n\rightarrow\infty}\frac{3\cdot5^n+2\cdot6^n}{2\cdot5^n+4\cdot7^n } \]}
{\( L=\lim\limits_{n\rightarrow\infty} \frac{3\cdot\left(\frac57\right)^n+2\cdot\left(\frac67\right)^n}{2\cdot\left(\frac57\right)^n+4} =0 \)}{\( L=\lim\limits_{n\rightarrow\infty}\frac{3\cdot\left(\frac56\right)^n+2}{2\cdot\left(\frac57\right)^n+4}=\frac12 \)}{\( L=\lim\limits_{n\rightarrow\infty}\frac{3+2\cdot\left(\frac65\right)^n}{2+4\cdot\left(\frac75\right)^n } =\frac32 \)}{\( L=\frac{3\cdot5^{\infty}+2\cdot6^{\infty}}{2\cdot5^{\infty}+4\cdot7^{\infty}}=\infty \)}{\( L=\lim\limits_{n\rightarrow\infty}\frac{3\cdot\left(\frac57\right)^n+2\cdot\left(\frac67\right)^n}{2\cdot\left(\frac57\right)^n+4\cdot\left(\frac77\right)^n}=\frac56 \)}{}}
\LANGsk{
\Question{
Vyberte správny výpočet limity postupnosti. \[ L=\lim\limits_{n\rightarrow\infty}\frac{3\cdot5^n+2\cdot6^n}{2\cdot5^n+4\cdot7^n } \]}
{\( L=\lim\limits_{n\rightarrow\infty} \frac{3\cdot\left(\frac57\right)^n+2\cdot\left(\frac67\right)^n}{2\cdot\left(\frac57\right)^n+4} =0 \)}{\( L=\lim\limits_{n\rightarrow\infty}\frac{3\cdot\left(\frac56\right)^n+2}{2\cdot\left(\frac57\right)^n+4}=\frac12 \)}{\( L=\lim\limits_{n\rightarrow\infty}\frac{3+2\cdot\left(\frac65\right)^n}{2+4\cdot\left(\frac75\right)^n } =\frac32 \)}{\( L=\frac{3\cdot5^{\infty}+2\cdot6^{\infty}}{2\cdot5^{\infty}+4\cdot7^{\infty}}=\infty \)}{\( L=\lim\limits_{n\rightarrow\infty}\frac{3\cdot\left(\frac57\right)^n+2\cdot\left(\frac67\right)^n}{2\cdot\left(\frac57\right)^n+4\cdot\left(\frac77\right)^n}=\frac56 \)}{}}
\LANGes{
\Question{
Elige la expresión correcta para calcular el límite de la siguiente sucesión: \[ L=\lim\limits_{n\rightarrow\infty}\frac{3\cdot5^n+2\cdot6^n}{2\cdot5^n+4\cdot7^n } \]}
{\( L=\lim\limits_{n\rightarrow\infty} \frac{3\cdot\left(\frac57\right)^n+2\cdot\left(\frac67\right)^n}{2\cdot\left(\frac57\right)^n+4} =0 \)}{\( L=\lim\limits_{n\rightarrow\infty}\frac{3\cdot\left(\frac56\right)^n+2}{2\cdot\left(\frac57\right)^n+4}=\frac12 \)}{\( L=\lim\limits_{n\rightarrow\infty}\frac{3+2\cdot\left(\frac65\right)^n}{2+4\cdot\left(\frac75\right)^n } =\frac32 \)}{\( L=\frac{3\cdot5^{\infty}+2\cdot6^{\infty}}{2\cdot5^{\infty}+4\cdot7^{\infty}}=\infty \)}{\( L=\lim\limits_{n\rightarrow\infty}\frac{3\cdot\left(\frac57\right)^n+2\cdot\left(\frac67\right)^n}{2\cdot\left(\frac57\right)^n+4\cdot\left(\frac77\right)^n}=\frac56 \)}{}}
\End

\Data\ID{1003047402}
\Updated{2021-05-03 08:05:07}
\Level{B}
\SubArea{Limit of a sequence}
\LANGen{
\Question{
Choose the best first step to simplify and calculate the limit of the following sequence. \[ \left(\frac{3\cdot5^n+2\cdot6^n}{2\cdot5^n+4\cdot6^n}\right)_{n=1}^{\infty} \]}
{We take \( 6^n \) outside the brackets in the numerator and the denominator.}{We take \( 5^n \) outside the brackets in the numerator and denominator.}{We divide the numerator and denominator by \( 5^n \).}{We divide the numerator by \( 6^n \).}{We divide the denominator by \( 6^n \).}{}}
\LANGpl{
\Question{
Wybierz pierwszy krok, by skutecznie uprościć i oszacować granicę ciągu.  \[ \left(\frac{3\cdot5^n+2\cdot6^n}{2\cdot5^n+4\cdot6^n}\right)_{n=1}^{\infty} \]}
{Wyciągamy  \( 6^n \) poza nawias w liczniku i mianowniku.}{Wyciągamy  \( 5^n \) poza nawias w liczniku i mianowniku.}{Dzielimy licznik i mianownik przez \( 5^n \).}{Dzielimy licznik przez \( 6^n \).}{Dzielimy mianownik przez \( 6^n \).}{}}
\LANGcs{
\Question{
Vyberte nejvhodnější první krok k úpravě a výpočtu limity následující posloupnosti. \[ \left(\frac{3\cdot5^n+2\cdot6^n}{2\cdot5^n+4\cdot6^n}\right)_{n=1}^{\infty} \]}
{Vytkneme v čitateli i jmenovateli \( 6^n \).}{Vytkneme v čitateli i jmenovateli \( 5^n \).}{Vydělíme čitatel i jmenovatel \( 5^n \).}{Vydělíme čitatel \( 6^n \).}{Vydělíme jmenovatel \( 6^n \).}{}}
\LANGsk{
\Question{
Vyberte najvhodnejší prvý krok k úprave a výpočtu limity nasledujúce postupnosti. \[ \left(\frac{3\cdot5^n+2\cdot6^n}{2\cdot5^n+4\cdot6^n}\right)_{n=1}^{\infty} \]}
{Vyjmeme v čitateli i menovateli \( 6^n \).}{Vyjmeme v čitateli i menovateli \( 5^n \).}{Vydelíme čitateľ i menovateľ \( 5^n \).}{Vydelíme čitateľ \( 6^n \).}{Vydelíme menovateľ \( 6^n \).}{}}
\LANGes{
\Question{
Elige el primer paso más adecuado para calcular el límite de la siguiente sucesión:
\[ \left(\frac{3\cdot5^n+2\cdot6^n}{2\cdot5^n+4\cdot6^n}\right)_{n=1}^{\infty} \]}
{Factorizar \( 6^n \) en el numerador y en el denominador.}{Factorizar \( 5^n \) en el numerador y en el denominador.}{Dividir el numerador y el denominador por \( 5^n \).}{Dividir el numerador por \( 6^n \).}{Dividir el denominador por \( 6^n \).}{}}
\End

\Data\ID{1003047403}
\Updated{2020-07-15 03:07:12}
\Level{B}
\SubArea{Limit of a sequence}
\LANGen{
\Question{
Find the limit \( \lim\limits_{n\rightarrow\infty}\frac{5^{n+1}+6^{n+1}}{5^{n+1}+6^n} \).}
{\( 6 \)}{\( 5 \)}{\( \frac56 \)}{\( 0 \)}{\( \infty \)}{}}
\LANGpl{
\Question{
Wyznacz granicę dla \( \lim\limits_{n\rightarrow\infty}\frac{5^{n+1}+6^{n+1}}{5^{n+1}+6^n} \).}
{\( 6 \)}{\( 5 \)}{\( \frac56 \)}{\( 0 \)}{\( \infty \)}{}}
\LANGcs{
\Question{
\( \lim\limits_{n\rightarrow\infty}\frac{5^{n+1}+6^{n+1}}{5^{n+1}+6^n} \) je rovna:}
{\( 6 \)}{\( 5 \)}{\( \frac56 \)}{\( 0 \)}{\( \infty \)}{}}
\LANGsk{
\Question{
\( \lim\limits_{n\rightarrow\infty}\frac{5^{n+1}+6^{n+1}}{5^{n+1}+6^n} \) je rovná:}
{\( 6 \)}{\( 5 \)}{\( \frac56 \)}{\( 0 \)}{\( \infty \)}{}}
\LANGes{
\Question{
Halla \( \lim\limits_{n\rightarrow\infty}\frac{5^{n+1}+6^{n+1}}{5^{n+1}+6^n} \).}
{\( 6 \)}{\( 5 \)}{\( \frac56 \)}{\( 0 \)}{\( \infty \)}{}}
\End

\Data\ID{1003047404}
\Updated{2020-07-15 03:07:12}
\Level{B}
\SubArea{Limit of a sequence}
\LANGen{
\Question{
Find the limit \( \lim\limits_{n\rightarrow\infty}\frac{5^{n+1}+6^n}{5^n+6^{n+1} } \).}
{\( \frac16 \)}{\( \frac56 \)}{\( 1 \)}{\( 0 \)}{\( \infty \)}{}}
\LANGpl{
\Question{
Wyznacz granicę dla  \( \lim\limits_{n\rightarrow\infty}\frac{5^{n+1}+6^n}{5^n+6^{n+1} } \).}
{\( \frac16 \)}{\( \frac56 \)}{\( 1 \)}{\( 0 \)}{\( \infty \)}{}}
\LANGcs{
\Question{
\( \lim\limits_{n\rightarrow\infty}\frac{5^{n+1}+6^n}{5^n+6^{n+1} } \) je rovna:}
{\( \frac16 \)}{\( \frac56 \)}{\( 1 \)}{\( 0 \)}{\( \infty \)}{}}
\LANGsk{
\Question{
\( \lim\limits_{n\rightarrow\infty}\frac{5^{n+1}+6^n}{5^n+6^{n+1} } \) je rovná:}
{\( \frac16 \)}{\( \frac56 \)}{\( 1 \)}{\( 0 \)}{\( \infty \)}{}}
\LANGes{
\Question{
Halla \( \lim\limits_{n\rightarrow\infty}\frac{5^{n+1}+6^n}{5^n+6^{n+1} } \).}
{\( \frac16 \)}{\( \frac56 \)}{\( 1 \)}{\( 0 \)}{\( \infty \)}{}}
\End

\Data\ID{1003047405}
\Updated{2020-07-15 03:07:12}
\Level{B}
\SubArea{Limit of a sequence}
\LANGen{
\Question{
The sequence  \( \left(\frac{3^n-4^{n-1}}{4^n}\right)_{n=1}^{\infty} \) is:}
{convergent and \(\lim\limits_{n\rightarrow\infty}\frac{3^n-4^{n-1}}{4^n}=-\frac14 \)}{convergent and \(\lim\limits_{n\rightarrow\infty}\frac{3^n-4^{n-1}}{4^n}=\frac14 \)}{convergent and \(\lim\limits_{n\rightarrow\infty}\frac{3^n-4^{n-1}}{4^n}=-1 \)}{convergent and \(\lim\limits_{n\rightarrow\infty}\frac{3^n-4^{n-1}}{4^n}=0 \)}{divergent}{}}
\LANGpl{
\Question{
Ciąg  \( \left(\frac{3^n-4^{n-1}}{4^n}\right)_{n=1}^{\infty} \) jest:}
{zbieżny i  \(\lim\limits_{n\rightarrow\infty}\frac{3^n-4^{n-1}}{4^n}=-\frac14 \)}{zbieżny i \(\lim\limits_{n\rightarrow\infty}\frac{3^n-4^{n-1}}{4^n}=\frac14 \)}{zbieżny i \(\lim\limits_{n\rightarrow\infty}\frac{3^n-4^{n-1}}{4^n}=-1 \)}{zbieżny i \(\lim\limits_{n\rightarrow\infty}\frac{3^n-4^{n-1}}{4^n}=0 \)}{rozbieżny}{}}
\LANGcs{
\Question{
Posloupnost  \( \left(\frac{3^n-4^{n-1}}{4^n}\right)_{n=1}^{\infty} \) je:}
{konvergentní a platí: \(\lim\limits_{n\rightarrow\infty}\frac{3^n-4^{n-1}}{4^n}=-\frac14 \)}{konvergentní a platí: \(\lim\limits_{n\rightarrow\infty}\frac{3^n-4^{n-1}}{4^n}=\frac14 \)}{konvergentní a platí: \(\lim\limits_{n\rightarrow\infty}\frac{3^n-4^{n-1}}{4^n}=-1 \)}{konvergentní a platí: \(\lim\limits_{n\rightarrow\infty}\frac{3^n-4^{n-1}}{4^n}=0 \)}{divergentní}{}}
\LANGsk{
\Question{
Postupnosť  \( \left(\frac{3^n-4^{n-1}}{4^n}\right)_{n=1}^{\infty} \) je:}
{konvergentná a platí: \(\lim\limits_{n\rightarrow\infty}\frac{3^n-4^{n-1}}{4^n}=-\frac14 \)}{konvergentná a platí: \(\lim\limits_{n\rightarrow\infty}\frac{3^n-4^{n-1}}{4^n}=\frac14 \)}{konvergentná a platí: \(\lim\limits_{n\rightarrow\infty}\frac{3^n-4^{n-1}}{4^n}=-1 \)}{konvergentná a platí: \(\lim\limits_{n\rightarrow\infty}\frac{3^n-4^{n-1}}{4^n}=0 \)}{divergentná}{}}
\LANGes{
\Question{
La sucesión  \( \left(\frac{3^n-4^{n-1}}{4^n}\right)_{n=1}^{\infty} \) es:}
{convergente y \(\lim\limits_{n\rightarrow\infty}\frac{3^n-4^{n-1}}{4^n}=-\frac14 \)}{convergente y \(\lim\limits_{n\rightarrow\infty}\frac{3^n-4^{n-1}}{4^n}=\frac14 \)}{convergente y \(\lim\limits_{n\rightarrow\infty}\frac{3^n-4^{n-1}}{4^n}=-1 \)}{convergente y \(\lim\limits_{n\rightarrow\infty}\frac{3^n-4^{n-1}}{4^n}=0 \)}{divergente}{}}
\End

\Data\ID{1003047406}
\Updated{2020-07-15 03:07:12}
\Level{B}
\SubArea{Limit of a sequence}
\LANGen{
\Question{
Choose the correct formula to calculate the limit. \[ L=\lim\limits_{n\rightarrow\infty}\frac{3^{n+1}+4^n}{2^n} \]}
{\( L=\lim\limits_{n\rightarrow\infty}\left(3\cdot\left(\frac32\right)^n+2^n\right)=\infty \)}{\( L=\lim\limits_{n\rightarrow\infty}\frac{2^n\left(3\cdot\left(\frac32\right)^n+2^n \right)}{2^n}=0 \)}{\(  L=\lim\limits_{n\rightarrow\infty}\frac{3^n \left(3+\left(\frac43\right)^n\right)}{2^n}=0 \)}{\(  L=\lim\limits_{n\rightarrow\infty}\frac{7^{n+1}}{2^n}=\infty \)}{\( L=\frac{3^{\infty+1}+4^{\infty}}{2^{\infty}} =\frac72 \)}{}}
\LANGpl{
\Question{
Wybierz odpowiedni wzór, za pomocą którego można obliczyć granicę ciągu. \[ L=\lim\limits_{n\rightarrow\infty}\frac{3^{n+1}+4^n}{2^n} \]}
{\( L=\lim\limits_{n\rightarrow\infty}\left(3\cdot\left(\frac32\right)^n+2^n\right)=\infty \)}{\( L=\lim\limits_{n\rightarrow\infty}\frac{2^n\left(3\cdot\left(\frac32\right)^n+2^n \right)}{2^n}=0 \)}{\(  L=\lim\limits_{n\rightarrow\infty}\frac{3^n \left(3+\left(\frac43\right)^n\right)}{2^n}=0 \)}{\(  L=\lim\limits_{n\rightarrow\infty}\frac{7^{n+1}}{2^n}=\infty \)}{\( L=\frac{3^{\infty+1}+4^{\infty}}{2^{\infty}} =\frac72 \)}{}}
\LANGcs{
\Question{
Vyberte správný výpočet limity posloupnosti. \[ L=\lim\limits_{n\rightarrow\infty}\frac{3^{n+1}+4^n}{2^n} \]}
{\( L=\lim\limits_{n\rightarrow\infty}\left(3\cdot\left(\frac32\right)^n+2^n\right)=\infty \)}{\( L=\lim\limits_{n\rightarrow\infty}\frac{2^n\left(3\cdot\left(\frac32\right)^n+2^n \right)}{2^n}=0 \)}{\(  L=\lim\limits_{n\rightarrow\infty}\frac{3^n \left(3+\left(\frac43\right)^n\right)}{2^n}=0 \)}{\(  L=\lim\limits_{n\rightarrow\infty}\frac{7^{n+1}}{2^n}=\infty \)}{\( L=\frac{3^{\infty+1}+4^{\infty}}{2^{\infty}} =\frac72 \)}{}}
\LANGsk{
\Question{
Vyberte správny výpočet limity postupnosti. \[ L=\lim\limits_{n\rightarrow\infty}\frac{3^{n+1}+4^n}{2^n} \]}
{\( L=\lim\limits_{n\rightarrow\infty}\left(3\cdot\left(\frac32\right)^n+2^n\right)=\infty \)}{\( L=\lim\limits_{n\rightarrow\infty}\frac{2^n\left(3\cdot\left(\frac32\right)^n+2^n \right)}{2^n}=0 \)}{\(  L=\lim\limits_{n\rightarrow\infty}\frac{3^n \left(3+\left(\frac43\right)^n\right)}{2^n}=0 \)}{\(  L=\lim\limits_{n\rightarrow\infty}\frac{7^{n+1}}{2^n}=\infty \)}{\( L=\frac{3^{\infty+1}+4^{\infty}}{2^{\infty}} =\frac72 \)}{}}
\LANGes{
\Question{
Elige la expresión correcta para calcular el límite de la siguiente sucesión: \[ L=\lim\limits_{n\rightarrow\infty}\frac{3^{n+1}+4^n}{2^n} \]}
{\( L=\lim\limits_{n\rightarrow\infty}\left(3\cdot\left(\frac32\right)^n+2^n\right)=\infty \)}{\( L=\lim\limits_{n\rightarrow\infty}\frac{2^n\left(3\cdot\left(\frac32\right)^n+2^n \right)}{2^n}=0 \)}{\(  L=\lim\limits_{n\rightarrow\infty}\frac{3^n \left(3+\left(\frac43\right)^n\right)}{2^n}=0 \)}{\(  L=\lim\limits_{n\rightarrow\infty}\frac{7^{n+1}}{2^n}=\infty \)}{\( L=\frac{3^{\infty+1}+4^{\infty}}{2^{\infty}} =\frac72 \)}{}}
\End

\Data\ID{1003047407}
\Updated{2020-07-15 03:07:12}
\Level{B}
\SubArea{Limit of a sequence}
\LANGen{
\Question{
Find the limit \( \lim\limits_{n\rightarrow\infty}\frac{3^{n+1}+4^{n+1}}{3^{n-1}+4^{n-1}}\).}
{\( 16 \)}{\( \frac1{16} \)}{\( 0 \)}{\( \infty \)}{\( 1 \)}{}}
\LANGpl{
\Question{
Wyznacz granicę dla  \( \lim\limits_{n\rightarrow\infty}\frac{3^{n+1}+4^{n+1}}{3^{n-1}+4^{n-1}}\).}
{\( 16 \)}{\( \frac1{16} \)}{\( 0 \)}{\( \infty \)}{\( 1 \)}{}}
\LANGcs{
\Question{
\( \lim\limits_{n\rightarrow\infty}\frac{3^{n+1}+4^{n+1}}{3^{n-1}+4^{n-1}}\) je rovna:}
{\( 16 \)}{\( \frac1{16} \)}{\( 0 \)}{\( \infty \)}{\( 1 \)}{}}
\LANGsk{
\Question{
\( \lim\limits_{n\rightarrow\infty}\frac{3^{n+1}+4^{n+1}}{3^{n-1}+4^{n-1}}\) je rovná:}
{\( 16 \)}{\( \frac1{16} \)}{\( 0 \)}{\( \infty \)}{\( 1 \)}{}}
\LANGes{
\Question{
Halla \( \lim\limits_{n\rightarrow\infty}\frac{3^{n+1}+4^{n+1}}{3^{n-1}+4^{n-1}}\).}
{\( 16 \)}{\( \frac1{16} \)}{\( 0 \)}{\( \infty \)}{\( 1 \)}{}}
\End

\Data\ID{1003047408}
\Updated{2021-05-03 08:05:25}
\Level{B}
\SubArea{Limit of a sequence}
\LANGen{
\Question{
Choose the best first step to simplify and calculate the limit \(\lim\limits_{n\rightarrow\infty}\frac{3^n+4^{n-1}}{3^n+4^{n+1}} \).}
{We divide the numerator and the denominator by \( 4^n \).}{We divide the numerator and denominator by \( 3^n \).}{We substitute \(n=\infty \).}{We take \( 3^n \) outside the brackets in the numerator and the denominator.}{We take \( 4 \) outside the brackets in the numerator and the denominator.}{}}
\LANGpl{
\Question{
Wybierz pierwszy krok by skutecznie uprościć i oszacować granicę \(\lim\limits_{n\rightarrow\infty}\frac{3^n+4^{n-1}}{3^n+4^{n+1}} \).}
{Dzielimy licznik i mianownik przez \( 4^n \).}{Dzielimy licznik i mianownik przez \( 3^n \).}{Podstawiamy  \(n=\infty \).}{Wyciągamy \( 3^n \) poza nawias w liczniku i mianowniku.}{Wyciągamy \( 4 \) poza nawias w liczniku i mianowniku.}{}}
\LANGcs{
\Question{
Vyberte nejvhodnější první krok k úpravě a výpočtu limity posloupnosti \(\lim\limits_{n\rightarrow\infty}\frac{3^n+4^{n-1}}{3^n+4^{n+1}} \).}
{Vydělíme čitatel i jmenovatel \( 4^n \).}{Vydělíme čitatel i jmenovatel \( 3^n \).}{Dosadíme \(n=\infty \).}{Vytkneme v čitateli i jmenovateli \( 3^n \).}{Vytkneme v čitateli i jmenovateli \( 4 \).}{}}
\LANGsk{
\Question{
Vyberte najvhodnejší prví krok k úprave a výpočtu limity postupnosti \(\lim\limits_{n\rightarrow\infty}\frac{3^n+4^{n-1}}{3^n+4^{n+1}} \).}
{Vydelíme čitateľ i menovateľ \( 4^n \).}{Vydelíme čitateľ i menovateľ \( 3^n \).}{Dosadíme \(n=\infty \).}{Vyjmeme v čitateli i menovateli \( 3^n \).}{Vyjmeme v čitateli i menovateli \( 4 \).}{}}
\LANGes{
\Question{
Elige el primer paso más adecuado para calcular el límite de la siguiente sucesión:
: \(\lim\limits_{n\rightarrow\infty}\frac{3^n+4^{n-1}}{3^n+4^{n+1}} \).}
{Dividir el numerador y el denominador por \( 4^n \).}{Dividir el numerador y el denominador por \( 3^n \).}{Sustituir \(n=\infty \).}{Factorizar \( 3^n \) en el numerador y en el denominador.}{Factorizar \( 4 \) en el numerador y en el denominador.}{}}
\End

\Data\ID{1003047409}
\Updated{2020-07-15 03:07:12}
\Level{B}
\SubArea{Limit of a sequence}
\LANGen{
\Question{
The sequence \( \left(\frac{2\cdot3^n+4^n+5}{4\cdot3^n-1}\right)_{n=1}^{\infty} \) is:}
{divergent and \( \lim\limits_{n\rightarrow\infty}\frac{2\cdot3^n+4^n+5}{4\cdot3^n-1}=\infty \)}{convergent and \( \lim\limits_{n\rightarrow\infty}\frac{2\cdot3^n+4^n+5}{4\cdot3^n-1}=\frac12 \)}{convergent and \( \lim\limits_{n\rightarrow\infty}\frac{2\cdot3^n+4^n+5}{4\cdot3^n-1}=\frac14 \)}{convergent and \( \lim\limits_{n\rightarrow\infty}\frac{2\cdot3^n+4^n+5}{4\cdot3^n-1}=0 \)}{divergent and it does not have an infinite limit}{}}
\LANGpl{
\Question{
Ciąg  \( \left(\frac{2\cdot3^n+4^n+5}{4\cdot3^n-1}\right)_{n=1}^{\infty} \) jest:}
{rozbieżny i \( \lim\limits_{n\rightarrow\infty}\frac{2\cdot3^n+4^n+5}{4\cdot3^n-1}=\infty \)}{zbieżny i \( \lim\limits_{n\rightarrow\infty}\frac{2\cdot3^n+4^n+5}{4\cdot3^n-1}=\frac12 \)}{zbieżny i \( \lim\limits_{n\rightarrow\infty}\frac{2\cdot3^n+4^n+5}{4\cdot3^n-1}=\frac14 \)}{zbieżny i \( \lim\limits_{n\rightarrow\infty}\frac{2\cdot3^n+4^n+5}{4\cdot3^n-1}=0 \)}{rozbieżny i nie ma nieskończonej granicy}{}}
\LANGcs{
\Question{
Posloupnost \( \left(\frac{2\cdot3^n+4^n+5}{4\cdot3^n-1}\right)_{n=1}^{\infty} \) je:}
{divergentní a platí: \( \lim\limits_{n\rightarrow\infty}\frac{2\cdot3^n+4^n+5}{4\cdot3^n-1}=\infty \)}{konvergentní a platí: \( \lim\limits_{n\rightarrow\infty}\frac{2\cdot3^n+4^n+5}{4\cdot3^n-1}=\frac12 \)}{konvergentní a platí: \( \lim\limits_{n\rightarrow\infty}\frac{2\cdot3^n+4^n+5}{4\cdot3^n-1}=\frac14 \)}{konvergentní a platí: \( \lim\limits_{n\rightarrow\infty}\frac{2\cdot3^n+4^n+5}{4\cdot3^n-1}=0 \)}{divergentní a nemá nevlastní limitu}{}}
\LANGsk{
\Question{
Postupnosť \( \left(\frac{2\cdot3^n+4^n+5}{4\cdot3^n-1}\right)_{n=1}^{\infty} \) je:}
{divergentná a platí: \( \lim\limits_{n\rightarrow\infty}\frac{2\cdot3^n+4^n+5}{4\cdot3^n-1}=\infty \)}{konvergentná a platí: \( \lim\limits_{n\rightarrow\infty}\frac{2\cdot3^n+4^n+5}{4\cdot3^n-1}=\frac12 \)}{konvergentná a platí: \( \lim\limits_{n\rightarrow\infty}\frac{2\cdot3^n+4^n+5}{4\cdot3^n-1}=\frac14 \)}{konvergentná a platí: \( \lim\limits_{n\rightarrow\infty}\frac{2\cdot3^n+4^n+5}{4\cdot3^n-1}=0 \)}{divergentná a nemá nevlastnú limitu}{}}
\LANGes{
\Question{
La sucesión \( \left(\frac{2\cdot3^n+4^n+5}{4\cdot3^n-1}\right)_{n=1}^{\infty} \) es:}
{divergente y \( \lim\limits_{n\rightarrow\infty}\frac{2\cdot3^n+4^n+5}{4\cdot3^n-1}=\infty \)}{convergente y \( \lim\limits_{n\rightarrow\infty}\frac{2\cdot3^n+4^n+5}{4\cdot3^n-1}=\frac12 \)}{convergente y \( \lim\limits_{n\rightarrow\infty}\frac{2\cdot3^n+4^n+5}{4\cdot3^n-1}=\frac14 \)}{convergente y \( \lim\limits_{n\rightarrow\infty}\frac{2\cdot3^n+4^n+5}{4\cdot3^n-1}=0 \)}{divergente y no tiene un límite infinito.}{}}
\End

\Data\ID{1003047410}
\Updated{2020-07-15 03:07:12}
\Level{B}
\SubArea{Limit of a sequence}
\LANGen{
\Question{
Find the limit \(\lim\limits_{n\rightarrow\infty}\frac{2\cdot3^n+4^n+5}{3\cdot5^n-1} \).}
{\( 0 \)}{\( \infty \)}{\( \frac23 \)}{\( -5 \)}{\( 4 \)}{}}
\LANGpl{
\Question{
Wyznacz granicę dla \(\lim\limits_{n\rightarrow\infty}\frac{2\cdot3^n+4^n+5}{3\cdot5^n-1} \).}
{\( 0 \)}{\( \infty \)}{\( \frac23 \)}{\( -5 \)}{\( 4 \)}{}}
\LANGcs{
\Question{
\(\lim\limits_{n\rightarrow\infty}\frac{2\cdot3^n+4^n+5}{3\cdot5^n-1} \) je rovna:}
{\( 0 \)}{\( \infty \)}{\( \frac23 \)}{\( -5 \)}{\( 4 \)}{}}
\LANGsk{
\Question{
\(\lim\limits_{n\rightarrow\infty}\frac{2\cdot3^n+4^n+5}{3\cdot5^n-1} \) je rovná:}
{\( 0 \)}{\( \infty \)}{\( \frac23 \)}{\( -5 \)}{\( 4 \)}{}}
\LANGes{
\Question{
Halla \(\lim\limits_{n\rightarrow\infty}\frac{2\cdot3^n+4^n+5}{3\cdot5^n-1} \).}
{\( 0 \)}{\( \infty \)}{\( \frac23 \)}{\( -5 \)}{\( 4 \)}{}}
\End

\Data\ID{1003047503}
\Updated{2020-07-15 03:07:12}
\Level{B}
\SubArea{Limit of a sequence}
\LANGen{
\Question{
Choose the sequence with the limit equal to \( -5 \).}
{\( (0.15^n-5)_{n=1}^{\infty} \)}{\( ((-5)^n-0.15)_{n=1}^{\infty} \)}{\( (5^n-5)_{n=1}^{\infty} \)}{\( (5^n+5)_{n=1}^{\infty} \)}{\( (-0.15^n+5)_{n=1}^{\infty} \)}{}}
\LANGpl{
\Question{
Wybierz ciąg, którego granica równa jest  \( -5 \).}
{\( (0{,}15^n-5)_{n=1}^{\infty} \)}{\( ((-5)^n-0{,}15)_{n=1}^{\infty} \)}{\( (5^n-5)_{n=1}^{\infty} \)}{\( (5^n+5)_{n=1}^{\infty} \)}{\( (-0{,}15^n+5)_{n=1}^{\infty} \)}{}}
\LANGcs{
\Question{
Vyberte posloupnost, jejíž limita je rovna \( -5 \).}
{\( (0{,}15^n-5)_{n=1}^{\infty} \)}{\( ((-5)^n-0{,}15)_{n=1}^{\infty} \)}{\( (5^n-5)_{n=1}^{\infty} \)}{\( (5^n+5)_{n=1}^{\infty} \)}{\( (-0{,}15^n+5)_{n=1}^{\infty} \)}{}}
\LANGsk{
\Question{
Vyberte postupnosť, ktorej limita je rovná \( -5 \).}
{\( (0{,}15^n-5)_{n=1}^{\infty} \)}{\( ((-5)^n-0{,}15)_{n=1}^{\infty} \)}{\( (5^n-5)_{n=1}^{\infty} \)}{\( (5^n+5)_{n=1}^{\infty} \)}{\( (-0{,}15^n+5)_{n=1}^{\infty} \)}{}}
\LANGes{
\Question{
Elige la sucesión cuyo límite es \( -5 \).}
{\( (0.15^n-5)_{n=1}^{\infty} \)}{\( ((-5)^n-0.15)_{n=1}^{\infty} \)}{\( (5^n-5)_{n=1}^{\infty} \)}{\( (5^n+5)_{n=1}^{\infty} \)}{\( (-0.15^n+5)_{n=1}^{\infty} \)}{}}
\End

\Data\ID{1003047506}
\Updated{2020-07-15 03:07:12}
\Level{B}
\SubArea{Limit of a sequence}
\LANGen{
\Question{
Choose the sequence with the limit equal to \(-5\).}
{\( \Bigl( \frac{(-3)^n+(-5)^{n+1}}{(-5)^n} \Bigr)_{n=1}^{\infty} \)}{\( \Bigl( \frac{(-3)^n-(-5)^{n+1}}{(-5)^n} \Bigr)_{n=1}^{\infty} \)}{\( \Bigl( \frac{(-7)^n+(-5)^{n+1}}{(-5)^n} \Bigr)_{n=1}^{\infty} \)}{\( \Bigl( \frac{(-3)^n+(-5)^{n+1}}{(-7)^n} \Bigr)_{n=1}^{\infty} \)}{\( \Bigl( \frac{(-5)^n+(-5)^{n+1}}{(-5)^n} \Bigr)_{n=1}^{\infty} \)}{}}
\LANGpl{
\Question{
Wybierz ciąg, którego granica równa jest \(-5\).}
{\( \Bigl( \frac{(-3)^n+(-5)^{n+1}}{(-5)^n} \Bigr)_{n=1}^{\infty} \)}{\( \Bigl( \frac{(-3)^n-(-5)^{n+1}}{(-5)^n} \Bigr)_{n=1}^{\infty} \)}{\( \Bigl( \frac{(-7)^n+(-5)^{n+1}}{(-5)^n} \Bigr)_{n=1}^{\infty} \)}{\( \Bigl( \frac{(-3)^n+(-5)^{n+1}}{(-7)^n} \Bigr)_{n=1}^{\infty} \)}{\( \Bigl( \frac{(-5)^n+(-5)^{n+1}}{(-5)^n} \Bigr)_{n=1}^{\infty} \)}{}}
\LANGcs{
\Question{
Vyberte posloupnost, jejíž limita je rovna \(-5\).}
{\( \Bigl( \frac{(-3)^n+(-5)^{n+1}}{(-5)^n} \Bigr)_{n=1}^{\infty} \)}{\( \Bigl( \frac{(-3)^n-(-5)^{n+1}}{(-5)^n} \Bigr)_{n=1}^{\infty} \)}{\( \Bigl( \frac{(-7)^n+(-5)^{n+1}}{(-5)^n} \Bigr)_{n=1}^{\infty} \)}{\( \Bigl( \frac{(-3)^n+(-5)^{n+1}}{(-7)^n} \Bigr)_{n=1}^{\infty} \)}{\( \Bigl( \frac{(-5)^n+(-5)^{n+1}}{(-5)^n} \Bigr)_{n=1}^{\infty} \)}{}}
\LANGsk{
\Question{
Vyberte postupnosť, ktorej limita je rovná \(-5\).}
{\( \Bigl( \frac{(-3)^n+(-5)^{n+1}}{(-5)^n} \Bigr)_{n=1}^{\infty} \)}{\( \Bigl( \frac{(-3)^n-(-5)^{n+1}}{(-5)^n} \Bigr)_{n=1}^{\infty} \)}{\( \Bigl( \frac{(-7)^n+(-5)^{n+1}}{(-5)^n} \Bigr)_{n=1}^{\infty} \)}{\( \Bigl( \frac{(-3)^n+(-5)^{n+1}}{(-7)^n} \Bigr)_{n=1}^{\infty} \)}{\( \Bigl( \frac{(-5)^n+(-5)^{n+1}}{(-5)^n} \Bigr)_{n=1}^{\infty} \)}{}}
\LANGes{
\Question{
Elige la sucesión cuyo límite es \(-5\).}
{\( \Bigl( \frac{(-3)^n+(-5)^{n+1}}{(-5)^n} \Bigr)_{n=1}^{\infty} \)}{\( \Bigl( \frac{(-3)^n-(-5)^{n+1}}{(-5)^n} \Bigr)_{n=1}^{\infty} \)}{\( \Bigl( \frac{(-7)^n+(-5)^{n+1}}{(-5)^n} \Bigr)_{n=1}^{\infty} \)}{\( \Bigl( \frac{(-3)^n+(-5)^{n+1}}{(-7)^n} \Bigr)_{n=1}^{\infty} \)}{\( \Bigl( \frac{(-5)^n+(-5)^{n+1}}{(-5)^n} \Bigr)_{n=1}^{\infty} \)}{}}
\End

\Data\ID{1003047507}
\Updated{2020-07-15 03:07:12}
\Level{B}
\SubArea{Limit of a sequence}
\LANGen{
\Question{
Find the following limit. \[ \lim\limits_{n\rightarrow\infty}\frac{4^{2n}}{4^n+2} \]}
{\( \infty \)}{\( 0 \)}{\( 1 \)}{\( 4 \)}{\( 16 \)}{}}
\LANGpl{
\Question{
Wyznacz granicę. \[ \lim\limits_{n\rightarrow\infty}\frac{4^{2n}}{4^n+2} \]}
{\( \infty \)}{\( 0 \)}{\( 1 \)}{\( 4 \)}{\( 16 \)}{}}
\LANGcs{
\Question{
\( \lim\limits_{n\rightarrow\infty}\frac{4^{2n}}{4^n+2} \) je rovna:}
{\( \infty \)}{\( 0 \)}{\( 1 \)}{\( 4 \)}{\( 16 \)}{}}
\LANGsk{
\Question{
\( \lim\limits_{n\rightarrow\infty}\frac{4^{2n}}{4^n+2} \) je rovná:}
{\( \infty \)}{\( 0 \)}{\( 1 \)}{\( 4 \)}{\( 16 \)}{}}
\LANGes{
\Question{
Halla: \[ \lim\limits_{n\rightarrow\infty}\frac{4^{2n}}{4^n+2} \]}
{\( \infty \)}{\( 0 \)}{\( 1 \)}{\( 4 \)}{\( 16 \)}{}}
\End

\Data\ID{1003047508}
\Updated{2020-07-15 03:07:12}
\Level{B}
\SubArea{Limit of a sequence}
\LANGen{
\Question{
Find the following limit. \[ \lim\limits_{n\rightarrow\infty}\frac{\left(\frac13\right)^n}{\left(\frac15\right)^n-\left(\frac14\right)^n} \]}
{\( -\infty \)}{\( \infty \)}{\( 0 \)}{\( 1 \)}{\( -1 \)}{}}
\LANGpl{
\Question{
Wyznacz granicę. \[ \lim\limits_{n\rightarrow\infty}\frac{\left(\frac13\right)^n}{\left(\frac15\right)^n-\left(\frac14\right)^n} \]}
{\( -\infty \)}{\( \infty \)}{\( 0 \)}{\( 1 \)}{\( -1 \)}{}}
\LANGcs{
\Question{
\( \lim\limits_{n\rightarrow\infty}\frac{\left(\frac13\right)^n}{\left(\frac15\right)^n-\left(\frac14\right)^n} \) je rovna:}
{\( -\infty \)}{\( \infty \)}{\( 0 \)}{\( 1 \)}{\( -1 \)}{}}
\LANGsk{
\Question{
\( \lim\limits_{n\rightarrow\infty}\frac{\left(\frac13\right)^n}{\left(\frac15\right)^n-\left(\frac14\right)^n} \) je rovná:}
{\( -\infty \)}{\( \infty \)}{\( 0 \)}{\( 1 \)}{\( -1 \)}{}}
\LANGes{
\Question{
Halla: \[ \lim\limits_{n\rightarrow\infty}\frac{\left(\frac13\right)^n}{\left(\frac15\right)^n-\left(\frac14\right)^n} \]}
{\( -\infty \)}{\( \infty \)}{\( 0 \)}{\( 1 \)}{\( -1 \)}{}}
\End

\Data\ID{1003047509}
\Updated{2020-07-15 03:07:12}
\Level{B}
\SubArea{Limit of a sequence}
\LANGen{
\Question{
Choose the sequence with the limit equal to \( -\frac25 \).}
{\( \left( \frac{2\log n-4}{3-5\log n}\right)_{n=1}^{\infty} \)}{\( \left( \frac{4\log n-2}{3-5\log n}\right)_{n=1}^{\infty} \)}{\( \left( \frac{2(\log n)^2-4}{3-5\log n}\right)_{n=1}^{\infty} \)}{\( \left( \frac{2\log n-4}{3-5(\log n)^2}\right)_{n=1}^{\infty} \)}{\( \left( \frac{4\log n-2}{3-5\log n}\right)_{n=1}^{\infty} \)}{}}
\LANGpl{
\Question{
Wybierz ciąg, którego granica równa jest \( -\frac25 \).}
{\( \left( \frac{2\log n-4}{3-5\log n}\right)_{n=1}^{\infty} \)}{\( \left( \frac{4\log n-2}{3-5\log n}\right)_{n=1}^{\infty} \)}{\( \left( \frac{2(\log n)^2-4}{3-5\log n}\right)_{n=1}^{\infty} \)}{\( \left( \frac{2\log n-4}{3-5(\log n)^2}\right)_{n=1}^{\infty} \)}{\( \left( \frac{4\log n-2}{3-5\log n}\right)_{n=1}^{\infty} \)}{}}
\LANGcs{
\Question{
Vyberte posloupnost, jejíž limita je rovna \( -\frac25 \).}
{\( \left( \frac{2\log n-4}{3-5\log n}\right)_{n=1}^{\infty} \)}{\( \left( \frac{4\log n-2}{3-5\log n}\right)_{n=1}^{\infty} \)}{\( \left( \frac{2(\log n)^2-4}{3-5\log n}\right)_{n=1}^{\infty} \)}{\( \left( \frac{2\log n-4}{3-5(\log n)^2}\right)_{n=1}^{\infty} \)}{\( \left( \frac{4\log n-2}{3-5\log n}\right)_{n=1}^{\infty} \)}{}}
\LANGsk{
\Question{
Vyberte postupnosť, ktorej limita je rovná \( -\frac25 \).}
{\( \left( \frac{2\log n-4}{3-5\log n}\right)_{n=1}^{\infty} \)}{\( \left( \frac{4\log n-2}{3-5\log n}\right)_{n=1}^{\infty} \)}{\( \left( \frac{2(\log n)^2-4}{3-5\log n}\right)_{n=1}^{\infty} \)}{\( \left( \frac{2\log n-4}{3-5(\log n)^2}\right)_{n=1}^{\infty} \)}{\( \left( \frac{4\log n-2}{3-5\log n}\right)_{n=1}^{\infty} \)}{}}
\LANGes{
\Question{
Elige la sucesión cuyo límite es \( -\frac25 \).}
{\( \left( \frac{2\log n-4}{3-5\log n}\right)_{n=1}^{\infty} \)}{\( \left( \frac{4\log n-2}{3-5\log n}\right)_{n=1}^{\infty} \)}{\( \left( \frac{2(\log n)^2-4}{3-5\log n}\right)_{n=1}^{\infty} \)}{\( \left( \frac{2\log n-4}{3-5(\log n)^2}\right)_{n=1}^{\infty} \)}{\( \left( \frac{4\log n-2}{3-5\log n}\right)_{n=1}^{\infty} \)}{}}
\End

\Data\ID{1003047510}
\Updated{2020-07-15 03:07:12}
\Level{B}
\SubArea{Limit of a sequence}
\LANGen{
\Question{
Choose the sequence with the limit equal to \( 0 \).}
{\( \left(\frac{3(\log n)^2+2\log n-1}{5(\log n)^3+2(\log n)^2+2}\right)_{n=1}^{\infty} \)}{\( \left(\frac{3(\log n)^3+2\log n-1}{5(\log n)^3+2(\log n)^2+2}\right)_{n=1}^{\infty} \)}{\( \left(\frac{3(\log n)^4+2\log n-1}{5(\log n)^3+2(\log n)^2+2}\right)_{n=1}^{\infty} \)}{\( \left(\frac{3(\log n)^3+2\log n-5}{5(\log n)^3-3(\log n)^2-2}\right)_{n=1}^{\infty} \)}{\( \left(\frac{3(\log n)^2+2\log n-1}{2(\log n)^2+2}\right)_{n=1}^{\infty} \)}{}}
\LANGpl{
\Question{
Wybierz ciąg, którego granica równa jest  \( 0 \).}
{\( \left(\frac{3(\log n)^2+2\log n-1}{5(\log n)^3+2(\log n)^2+2}\right)_{n=1}^{\infty} \)}{\( \left(\frac{3(\log n)^3+2\log n-1}{5(\log n)^3+2(\log n)^2+2}\right)_{n=1}^{\infty} \)}{\( \left(\frac{3(\log n)^4+2\log n-1}{5(\log n)^3+2(\log n)^2+2}\right)_{n=1}^{\infty} \)}{\( \left(\frac{3(\log n)^3+2\log n-5}{5(\log n)^3-3(\log n)^2-2}\right)_{n=1}^{\infty} \)}{\( \left(\frac{3(\log n)^2+2\log n-1}{2(\log n)^2+2}\right)_{n=1}^{\infty} \)}{}}
\LANGcs{
\Question{
Vyberte posloupnost, jejíž limita je rovna \( 0 \).}
{\( \left(\frac{3(\log n)^2+2\log n-1}{5(\log n)^3+2(\log n)^2+2}\right)_{n=1}^{\infty} \)}{\( \left(\frac{3(\log n)^3+2\log n-1}{5(\log n)^3+2(\log n)^2+2}\right)_{n=1}^{\infty} \)}{\( \left(\frac{3(\log n)^4+2\log n-1}{5(\log n)^3+2(\log n)^2+2}\right)_{n=1}^{\infty} \)}{\( \left(\frac{3(\log n)^3+2\log n-5}{5(\log n)^3-3(\log n)^2-2}\right)_{n=1}^{\infty} \)}{\( \left(\frac{3(\log n)^2+2\log n-1}{2(\log n)^2+2}\right)_{n=1}^{\infty} \)}{}}
\LANGsk{
\Question{
Vyberte postupnosť, ktorej limita je rovná \( 0 \).}
{\( \left(\frac{3(\log n)^2+2\log n-1}{5(\log n)^3+2(\log n)^2+2}\right)_{n=1}^{\infty} \)}{\( \left(\frac{3(\log n)^3+2\log n-1}{5(\log n)^3+2(\log n)^2+2}\right)_{n=1}^{\infty} \)}{\( \left(\frac{3(\log n)^4+2\log n-1}{5(\log n)^3+2(\log n)^2+2}\right)_{n=1}^{\infty} \)}{\( \left(\frac{3(\log n)^3+2\log n-5}{5(\log n)^3-3(\log n)^2-2}\right)_{n=1}^{\infty} \)}{\( \left(\frac{3(\log n)^2+2\log n-1}{2(\log n)^2+2}\right)_{n=1}^{\infty} \)}{}}
\LANGes{
\Question{
Elige la sucesión cuyo límite es \( 0 \).}
{\( \left(\frac{3(\log n)^2+2\log n-1}{5(\log n)^3+2(\log n)^2+2}\right)_{n=1}^{\infty} \)}{\( \left(\frac{3(\log n)^3+2\log n-1}{5(\log n)^3+2(\log n)^2+2}\right)_{n=1}^{\infty} \)}{\( \left(\frac{3(\log n)^4+2\log n-1}{5(\log n)^3+2(\log n)^2+2}\right)_{n=1}^{\infty} \)}{\( \left(\frac{3(\log n)^3+2\log n-5}{5(\log n)^3-3(\log n)^2-2}\right)_{n=1}^{\infty} \)}{\( \left(\frac{3(\log n)^2+2\log n-1}{2(\log n)^2+2}\right)_{n=1}^{\infty} \)}{}}
\End

\Data\ID{9000063602}
\Updated{2020-07-15 03:07:37}
\Level{B}
\SubArea{Limit of a sequence}
\LANGen{
\Question{
Find the following limit. 
\[ 
 \lim _{n\to \infty }(-1)^{n} \frac{3} 
{2n + 1}
\]}
{\(0\)}{\(-\frac{3} 
{2}\)}{\(\frac{3}
{2}\)}{\(- 1\)}{}{}}
\LANGpl{
\Question{
Wyznacz granicę ciągu.
\[ 
 \lim _{n\to \infty }(-1)^{n} \frac{3} 
{2n + 1}
\]}
{\(0\)}{\(-\frac{3} 
{2}\)}{\(\frac{3}
{2}\)}{\(- 1\)}{}{}}
\LANGcs{
\Question{
\(\lim _{n\to \infty }(-1)^{n} \frac{3} 
{2n+1}\) je
rovna:}
{\(0\)}{\(-\frac{3} 
{2}\)}{\(\frac{3}
{2}\)}{\(- 1\)}{}{}}
\LANGsk{
\Question{
\(\lim _{n\to \infty }(-1)^{n} \frac{3} 
{2n+1}\) je
rovná:}
{\(0\)}{\(-\frac{3} 
{2}\)}{\(\frac{3}
{2}\)}{\(- 1\)}{}{}}
\LANGes{
\Question{
Halla:
\[ 
 \lim _{n\to \infty }(-1)^{n} \frac{3} 
{2n + 1}
\]}
{\(0\)}{\(-\frac{3} 
{2}\)}{\(\frac{3}
{2}\)}{\(- 1\)}{}{}}
\End

\Data\ID{9000063604}
\Updated{2020-07-15 03:07:37}
\Level{B}
\SubArea{Limit of a sequence}
\LANGen{
\Question{
Find the following limit.
\[ 
 \lim _{n\to \infty }\sin 2\pi n
\]}
{\(0\)}{\(1\)}{\(- 1\)}{\(\infty \)}{}{}}
\LANGpl{
\Question{
Wyznacz granicę ciągu.
\[ 
 \lim _{n\to \infty }\sin 2\pi n
\]}
{\(0\)}{\(1\)}{\(- 1\)}{\(\infty \)}{}{}}
\LANGcs{
\Question{
\(\lim _{n\to \infty }\sin 2\pi n\) je
rovna:}
{\(0\)}{\(1\)}{\(- 1\)}{\(\infty \)}{}{}}
\LANGsk{
\Question{
\(\lim _{n\to \infty }\sin 2\pi n\) je
rovná:}
{\(0\)}{\(1\)}{\(- 1\)}{\(\infty \)}{}{}}
\LANGes{
\Question{
Halla:
\[ 
 \lim _{n\to \infty }\sin 2\pi n
\]}
{\(0\)}{\(1\)}{\(- 1\)}{\(\infty \)}{}{}}
\End

\Data\ID{9000063605}
\Updated{2020-07-15 03:07:37}
\Level{B}
\SubArea{Limit of a sequence}
\LANGen{
\Question{
Find the following limit.
\[ 
 \lim _{n\to \infty }\log 3^{n}
\]}
{\(\infty \)}{\(- 1\)}{\(0\)}{\(3\)}{}{}}
\LANGpl{
\Question{
Wyznacz granicę ciągu.
\[ 
 \lim _{n\to \infty }\log 3^{n}
\]}
{\(\infty \)}{\(- 1\)}{\(0\)}{\(3\)}{}{}}
\LANGcs{
\Question{
\(\lim _{n\to \infty }\log 3^{n}\) je
rovna:}
{\(\infty \)}{\(- 1\)}{\(0\)}{\(3\)}{}{}}
\LANGsk{
\Question{
\(\lim _{n\to \infty }\log 3^{n}\) je
rovná:}
{\(\infty \)}{\(- 1\)}{\(0\)}{\(3\)}{}{}}
\LANGes{
\Question{
Halla:
\[ 
 \lim _{n\to \infty }\log 3^{n}
\]}
{\(\infty \)}{\(- 1\)}{\(0\)}{\(3\)}{}{}}
\End

\Data\ID{9000063607}
\Updated{2020-07-15 03:07:37}
\Level{B}
\SubArea{Limit of a sequence}
\LANGen{
\Question{
Find the following limit.
\[ 
 \lim _{n\to \infty } \frac{1} 
{\log 10^{n}}
\]}
{\(0\)}{\(1\)}{\(10\)}{\(\infty \)}{}{}}
\LANGpl{
\Question{
Wyznacz granicę ciągu.
\[ 
 \lim _{n\to \infty } \frac{1} 
{\log 10^{n}}
\]}
{\(0\)}{\(1\)}{\(10\)}{\(\infty \)}{}{}}
\LANGcs{
\Question{
\(\lim _{n\to \infty } \frac{1} 
{\log 10^{n}}\) je
rovna:}
{\(0\)}{\(1\)}{\(10\)}{\(\infty \)}{}{}}
\LANGsk{
\Question{
\(\lim _{n\to \infty } \frac{1} 
{\log 10^{n}}\) je
rovná:}
{\(0\)}{\(1\)}{\(10\)}{\(\infty \)}{}{}}
\LANGes{
\Question{
Halla:
\[ 
 \lim _{n\to \infty } \frac{1} 
{\log 10^{n}}
\]}
{\(0\)}{\(1\)}{\(10\)}{\(\infty \)}{}{}}
\End

\Data\ID{9000063608}
\Updated{2020-07-15 03:07:37}
\Level{B}
\SubArea{Limit of a sequence}
\LANGen{
\Question{
Find the following limit.
\[ 
 \lim _{n\to \infty }\frac{2^{n} + 3^{n}} 
 {3^{n}}
\]}
{\(1\)}{\(2\)}{\(3\)}{\(\infty \)}{}{}}
\LANGpl{
\Question{
Wyznacz granicę ciągu.
\[ 
 \lim _{n\to \infty }\frac{2^{n} + 3^{n}} 
 {3^{n}}
\]}
{\(1\)}{\(2\)}{\(3\)}{\(\infty \)}{}{}}
\LANGcs{
\Question{
\(\lim _{n\to \infty }\frac{2^{n}+3^{n}} 
 {3^{n}} \) je
rovna:}
{\(1\)}{\(2\)}{\(3\)}{\(\infty \)}{}{}}
\LANGsk{
\Question{
\(\lim _{n\to \infty }\frac{2^{n}+3^{n}} 
 {3^{n}} \) je
rovná:}
{\(1\)}{\(2\)}{\(3\)}{\(\infty \)}{}{}}
\LANGes{
\Question{
Halla:
\[ 
 \lim _{n\to \infty }\frac{2^{n} + 3^{n}} 
 {3^{n}}
\]}
{\(1\)}{\(2\)}{\(3\)}{\(\infty \)}{}{}}
\End

\Data\ID{1003047501}
\Updated{2020-07-15 03:07:12}
\Level{C}
\SubArea{Limit of a sequence}
\LANGen{
\Question{
Find the following limit. \[ \lim\limits_{n\rightarrow\infty}(2\sqrt{n}-3) \]}
{\( \infty \)}{\( 2 \)}{\( -1 \)}{\( 0 \)}{\( -3 \)}{}}
\LANGpl{
\Question{
Wyznacz granicę. \[ \lim\limits_{n\rightarrow\infty}(2\sqrt{n}-3) \]}
{\( \infty \)}{\( 2 \)}{\( -1 \)}{\( 0 \)}{\( -3 \)}{}}
\LANGcs{
\Question{
\( \lim\limits_{n\rightarrow\infty}(2\sqrt{n}-3) \) je rovna:}
{\( \infty \)}{\( 2 \)}{\( -1 \)}{\( 0 \)}{\( -3 \)}{}}
\LANGsk{
\Question{
\( \lim\limits_{n\rightarrow\infty}(2\sqrt{n}-3) \) je rovná:}
{\( \infty \)}{\( 2 \)}{\( -1 \)}{\( 0 \)}{\( -3 \)}{}}
\LANGes{
\Question{
Halla: \[ \lim\limits_{n\rightarrow\infty}(2\sqrt{n}-3) \]}
{\( \infty \)}{\( 2 \)}{\( -1 \)}{\( 0 \)}{\( -3 \)}{}}
\End

\Data\ID{1003047502}
\Updated{2020-07-15 03:07:12}
\Level{C}
\SubArea{Limit of a sequence}
\LANGen{
\Question{
Find the following limit. \[ \lim\limits_{n\rightarrow\infty}\frac 4{2\sqrt{n}-3} \]}
{\( 0 \)}{\( 2 \)}{\( \frac12 \)}{\( 4 \)}{\( \infty \)}{}}
\LANGpl{
\Question{
Wyznacz granicę. \[ \lim\limits_{n\rightarrow\infty}\frac 4{2\sqrt{n}-3} \]}
{\( 0 \)}{\( 2 \)}{\( \frac12 \)}{\( 4 \)}{\( \infty \)}{}}
\LANGcs{
\Question{
\( \lim\limits_{n\rightarrow\infty}\frac 4{2\sqrt{n}-3} \) je rovna:}
{\( 0 \)}{\( 2 \)}{\( \frac12 \)}{\( 4 \)}{\( \infty \)}{}}
\LANGsk{
\Question{
\( \lim\limits_{n\rightarrow\infty}\frac 4{2\sqrt{n}-3} \) je rovná:}
{\( 0 \)}{\( 2 \)}{\( \frac12 \)}{\( 4 \)}{\( \infty \)}{}}
\LANGes{
\Question{
Halla: \[ \lim\limits_{n\rightarrow\infty}\frac 4{2\sqrt{n}-3} \]}
{\( 0 \)}{\( 2 \)}{\( \frac12 \)}{\( 4 \)}{\( \infty \)}{}}
\End

\Data\ID{1003047505}
\Updated{2020-07-15 03:07:12}
\Level{C}
\SubArea{Limit of a sequence}
\LANGen{
\Question{
Find the following limit. \[ \lim\limits_{n\rightarrow\infty}\Bigl(\frac5{\sqrt{n}}-\frac2{3^n} +\frac{(-1)^n}{2n^2-1}-7\Bigr) \]}
{\( -7 \)}{\( 0 \)}{\( -4 \)}{\( \infty \)}{\( -\infty \)}{}}
\LANGpl{
\Question{
Wyznacz granicę. \[ \lim\limits_{n\rightarrow\infty}\Bigl(\frac5{\sqrt{n}}-\frac2{3^n} +\frac{(-1)^n}{2n^2-1}-7\Bigr) \]}
{\( -7 \)}{\( 0 \)}{\( -4 \)}{\( \infty \)}{\( -\infty \)}{}}
\LANGcs{
\Question{
\( \lim\limits_{n\rightarrow\infty}\Bigl(\frac5{\sqrt{n}}-\frac2{3^n} +\frac{(-1)^n}{2n^2-1}-7\Bigr) \) je rovna:}
{\( -7 \)}{\( 0 \)}{\( -4 \)}{\( \infty \)}{\( -\infty \)}{}}
\LANGsk{
\Question{
\( \lim\limits_{n\rightarrow\infty}\Bigl(\frac5{\sqrt{n}}-\frac2{3^n} +\frac{(-1)^n}{2n^2-1}-7\Bigr) \) je rovná:}
{\( -7 \)}{\( 0 \)}{\( -4 \)}{\( \infty \)}{\( -\infty \)}{}}
\LANGes{
\Question{
Halla: \[ \lim\limits_{n\rightarrow\infty}\Bigl(\frac5{\sqrt{n}}-\frac2{3^n} +\frac{(-1)^n}{2n^2-1}-7\Bigr) \]}
{\( -7 \)}{\( 0 \)}{\( -4 \)}{\( \infty \)}{\( -\infty \)}{}}
\End

\Data\ID{1003047601}
\Updated{2020-07-15 03:07:12}
\Level{C}
\SubArea{Limit of a sequence}
\LANGen{
\Question{
Find the  limit  \( \lim\limits_{n\to\infty}\left(n-\sqrt{n-1}\right) \).}
{\( \infty \)}{\( 0 \)}{\( -\infty \)}{\( 1 \)}{\( \frac12 \)}{}}
\LANGpl{
\Question{
Wyznacz granicę dla  \( \lim\limits_{n\to\infty}\left(n-\sqrt{n-1}\right) \).}
{\( \infty \)}{\( 0 \)}{\( -\infty \)}{\( 1 \)}{\( \frac12 \)}{}}
\LANGcs{
\Question{
\( \lim\limits_{n\to\infty}\left(n-\sqrt{n-1}\right) \) je rovna:}
{\( \infty \)}{\( 0 \)}{\( -\infty \)}{\( 1 \)}{\( \frac12 \)}{}}
\LANGsk{
\Question{
\( \lim\limits_{n\to\infty}\left(n-\sqrt{n-1}\right) \) je rovná:}
{\( \infty \)}{\( 0 \)}{\( -\infty \)}{\( 1 \)}{\( \frac12 \)}{}}
\LANGes{
\Question{
Halla  \( \lim\limits_{n\to\infty}\left(n-\sqrt{n-1}\right) \).}
{\( \infty \)}{\( 0 \)}{\( -\infty \)}{\( 1 \)}{\( \frac12 \)}{}}
\End

\Data\ID{1003047602}
\Updated{2021-05-03 08:05:50}
\Level{C}
\SubArea{Limit of a sequence}
\LANGen{
\Question{
Choose the step to take first to efficiently evaluate the limit of the sequence \( \left(n-\sqrt{n^2-1} \right)_{n=1}^{\infty} \).}
{We expand with  the expression \( n+\sqrt{n^2-1} \).}{We expand with  the expression  \( n-\sqrt{n^2-1} \).}{We expand with  \( n \).}{We multiply by the expression \( n+\sqrt{n^2-1} \).}{We multiply by the expression \( n-\sqrt{n^2-1} \).}{We substitute \( n=\infty \).}}
\LANGpl{
\Question{
Wybierz pierwszy krok do skutecznego obliczenia granicy ciągu \( \left(n-\sqrt{n^2-1} \right)_{n=1}^{\infty} \).}
{Rozszerzamy  z wyrażeniem \( n+\sqrt{n^2-1} \).}{Rozszerzamy z wyrażeniem \( n-\sqrt{n^2-1} \).}{Rozszerzamy z   \( n \).}{Mnożymy przez wyrażenie \( n+\sqrt{n^2-1} \).}{Mnożymy przez wyrażenie \( n-\sqrt{n^2-1} \).}{Dzielimy przez \( n=\infty \).}}
\LANGcs{
\Question{
Vyberte vhodný postup pro výpočet limity posloupnosti \( \left(n-\sqrt{n^2-1} \right)_{n=1}^{\infty} \).}
{Rozšíříme výrazem  \( n+\sqrt{n^2-1} \).}{Rozšíříme výrazem  \( n-\sqrt{n^2-1} \).}{Rozšíříme  \( n \).}{Vynásobíme výrazem \( n+\sqrt{n^2-1} \).}{Vynásobíme výrazem \( n-\sqrt{n^2-1} \).}{Dosadíme \( n=\infty \).}}
\LANGsk{
\Question{
Vyberte vhodný postup pre výpočet limity postupnosti \( \left(n-\sqrt{n^2-1} \right)_{n=1}^{\infty} \).}
{Rozšírime výrazom  \( n+\sqrt{n^2-1} \).}{Rozšírime výrazom  \( n-\sqrt{n^2-1} \).}{Rozšírime  \( n \).}{Vynásobíme výrazom \( n+\sqrt{n^2-1} \).}{Vynásobíme výrazom \( n-\sqrt{n^2-1} \).}{Dosadíme \( n=\infty \).}}
\LANGes{
\Question{
Elige el primer paso más adecuado para calcular el límite de la siguiente sucesión:
\( \left(n-\sqrt{n^2-1} \right)_{n=1}^{\infty} \).}
{Expandir con la expresión \( n+\sqrt{n^2-1} \).}{Expandir con la expresión \( n-\sqrt{n^2-1} \).}{Expandir con la expresión \( n \).}{Multiplicar por la expresión \( n+\sqrt{n^2-1} \).}{Multiplicar por la expresión \( n-\sqrt{n^2-1} \).}{Sustituir \( n=\infty \).}}
\End

\Data\ID{1003047603}
\Updated{2020-07-15 03:07:12}
\Level{C}
\SubArea{Limit of a sequence}
\LANGen{
\Question{
Find the  limit  \( \lim\limits_{n\to\infty}\left( \sqrt{4n^2+3n}-2n \right) \).}
{\( \frac34 \)}{\( \infty \)}{\( 0 \)}{\( -\infty \)}{\( \sqrt2 \)}{}}
\LANGpl{
\Question{
Wyznacz granicę dla \( \lim\limits_{n\to\infty}\left( \sqrt{4n^2+3n}-2n \right) \).}
{\( \frac34 \)}{\( \infty \)}{\( 0 \)}{\( -\infty \)}{\( \sqrt2 \)}{}}
\LANGcs{
\Question{
\( \lim\limits_{n\to\infty}\left( \sqrt{4n^2+3n}-2n \right) \) je rovna:}
{\( \frac34 \)}{\( \infty \)}{\( 0 \)}{\( -\infty \)}{\( \sqrt2 \)}{}}
\LANGsk{
\Question{
\( \lim\limits_{n\to\infty}\left( \sqrt{4n^2+3n}-2n \right) \) je rovná:}
{\( \frac34 \)}{\( \infty \)}{\( 0 \)}{\( -\infty \)}{\( \sqrt2 \)}{}}
\LANGes{
\Question{
Halla:  \( \lim\limits_{n\to\infty}\left( \sqrt{4n^2+3n}-2n \right) \).}
{\( \frac34 \)}{\( \infty \)}{\( 0 \)}{\( -\infty \)}{\( \sqrt2 \)}{}}
\End

\Data\ID{1003047604}
\Updated{2020-07-15 03:07:12}
\Level{C}
\SubArea{Limit of a sequence}
\LANGen{
\Question{
Choose the correct computation  of the limit.  \[ L=\lim\limits_{n\to\infty} \left( \sqrt{n^2+3n}-2n \right) \]}
{\( L=\lim\limits_{n\to\infty}n\left( \sqrt{1+\frac3n}-2 \right) = -\infty \)}{\( L= \infty-\infty=0 \)}{\( L=\lim\limits_{n\to\infty}(n-2n)=-\infty \)}{\( L=\lim\limits_{n\to\infty} \left( n^2+3n-4n^2 \right) =-3 \)}{\( L=\lim\limits_{n\to\infty}\frac{n^2+3n-4n^2}{\sqrt{n^2+3n}+2n}=\infty \)}{}}
\LANGpl{
\Question{
Wybierz poprawne obliczenie granicy. \[ L=\lim\limits_{n\to\infty} \left( \sqrt{n^2+3n}-2n \right) \]}
{\( L=\lim\limits_{n\to\infty}n\left( \sqrt{1+\frac3n}-2 \right) = -\infty \)}{\( L= \infty-\infty=0 \)}{\( L=\lim\limits_{n\to\infty}(n-2n)=-\infty \)}{\( L=\lim\limits_{n\to\infty} \left( n^2+3n-4n^2 \right) =-3 \)}{\( L=\lim\limits_{n\to\infty}\frac{n^2+3n-4n^2}{\sqrt{n^2+3n}+2n}=\infty \)}{}}
\LANGcs{
\Question{
Vyberte správný výpočet limity posloupnosti.  \[ L=\lim\limits_{n\to\infty} \left( \sqrt{n^2+3n}-2n \right) \]}
{\( L=\lim\limits_{n\to\infty}n\left( \sqrt{1+\frac3n}-2 \right) = -\infty \)}{\( L= \infty-\infty=0 \)}{\( L=\lim\limits_{n\to\infty}(n-2n)=-\infty \)}{\( L=\lim\limits_{n\to\infty} \left( n^2+3n-4n^2 \right) =-3 \)}{\( L=\lim\limits_{n\to\infty}\frac{n^2+3n-4n^2}{\sqrt{n^2+3n}+2n}=\infty \)}{}}
\LANGsk{
\Question{
Vyberte správny výpočet limity postupnosti.  \[ L=\lim\limits_{n\to\infty} \left( \sqrt{n^2+3n}-2n \right) \]}
{\( L=\lim\limits_{n\to\infty}n\left( \sqrt{1+\frac3n}-2 \right) = -\infty \)}{\( L= \infty-\infty=0 \)}{\( L=\lim\limits_{n\to\infty}(n-2n)=-\infty \)}{\( L=\lim\limits_{n\to\infty} \left( n^2+3n-4n^2 \right) =-3 \)}{\( L=\lim\limits_{n\to\infty}\frac{n^2+3n-4n^2}{\sqrt{n^2+3n}+2n}=\infty \)}{}}
\LANGes{
\Question{
Elige la expresión correcta para calcular el límite de la siguiente sucesión:  \[ L=\lim\limits_{n\to\infty} \left( \sqrt{n^2+3n}-2n \right) \]}
{\( L=\lim\limits_{n\to\infty}n\left( \sqrt{1+\frac3n}-2 \right) = -\infty \)}{\( L= \infty-\infty=0 \)}{\( L=\lim\limits_{n\to\infty}(n-2n)=-\infty \)}{\( L=\lim\limits_{n\to\infty} \left( n^2+3n-4n^2 \right) =-3 \)}{\( L=\lim\limits_{n\to\infty}\frac{n^2+3n-4n^2}{\sqrt{n^2+3n}+2n}=\infty \)}{}}
\End

\Data\ID{1003047605}
\Updated{2020-07-15 03:07:12}
\Level{C}
\SubArea{Limit of a sequence}
\LANGen{
\Question{
Find the limit \( \lim\limits_{n\to\infty} \left( \sqrt n-\sqrt{n-1} \right) \).}
{\( 0 \)}{\( \infty \)}{\( -\infty \)}{\( -1 \)}{\( \sqrt2 \)}{}}
\LANGpl{
\Question{
Wyznacz granicę dla \( \lim\limits_{n\to\infty} \left( \sqrt n-\sqrt{n-1} \right) \).}
{\( 0 \)}{\( \infty \)}{\( -\infty \)}{\( -1 \)}{\( \sqrt2 \)}{}}
\LANGcs{
\Question{
\( \lim\limits_{n\to\infty} \left( \sqrt n-\sqrt{n-1} \right) \) je rovna:}
{\( 0 \)}{\( \infty \)}{\( -\infty \)}{\( -1 \)}{\( \sqrt2 \)}{}}
\LANGsk{
\Question{
\( \lim\limits_{n\to\infty} \left( \sqrt n-\sqrt{n-1} \right) \) je rovná:}
{\( 0 \)}{\( \infty \)}{\( -\infty \)}{\( -1 \)}{\( \sqrt2 \)}{}}
\LANGes{
\Question{
Halla \( \lim\limits_{n\to\infty} \left( \sqrt n-\sqrt{n-1} \right) \).}
{\( 0 \)}{\( \infty \)}{\( -\infty \)}{\( -1 \)}{\( \sqrt2 \)}{}}
\End

\Data\ID{1003047606}
\Updated{2020-07-15 03:07:12}
\Level{C}
\SubArea{Limit of a sequence}
\LANGen{
\Question{
The sequence \( \left( \sqrt n \left( \sqrt n-\sqrt{n-1} \right) \right)_{n=1}^{\infty} \) is:}
{convergent and \( \lim\limits_{n\to\infty} \sqrt n \left( \sqrt n-\sqrt{n-1} \right) =\frac12 \)}{convergent and \( \lim\limits_{n\to\infty} \sqrt n \left( \sqrt n-\sqrt{n-1} \right) =0 \)}{convergent and \( \lim\limits_{n\to\infty} \sqrt n \left( \sqrt n-\sqrt{n-1} \right) =2 \)}{divergent and \( \lim\limits_{n\to\infty} \sqrt n \left( \sqrt n-\sqrt{n-1} \right) =\infty \)}{divergent and it does not have an infinite limit}{}}
\LANGpl{
\Question{
Ciąg \( \left( \sqrt n \left( \sqrt n-\sqrt{n-1} \right) \right)_{n=1}^{\infty} \) jest:}
{zbieżny i \( \lim\limits_{n\to\infty} \sqrt n \left( \sqrt n-\sqrt{n-1} \right) =\frac12 \)}{zbieżny i \( \lim\limits_{n\to\infty} \sqrt n \left( \sqrt n-\sqrt{n-1} \right) =0 \)}{zbieżny i \( \lim\limits_{n\to\infty} \sqrt n \left( \sqrt n-\sqrt{n-1} \right) =2 \)}{rozbieżny i \( \lim\limits_{n\to\infty} \sqrt n \left( \sqrt n-\sqrt{n-1} \right) =\infty \)}{rozbieżny i nie ma nieskończonej granicy}{}}
\LANGcs{
\Question{
Posloupnost \( \left( \sqrt n \left( \sqrt n-\sqrt{n-1} \right) \right)_{n=1}^{\infty} \) je:}
{konvergentní a platí: \( \lim\limits_{n\to\infty} \sqrt n \left( \sqrt n-\sqrt{n-1} \right) =\frac12 \)}{konvergentní a platí: \( \lim\limits_{n\to\infty} \sqrt n \left( \sqrt n-\sqrt{n-1} \right) =0 \)}{konvergentní a platí: \( \lim\limits_{n\to\infty} \sqrt n \left( \sqrt n-\sqrt{n-1} \right) =2 \)}{divergentní a platí: \( \lim\limits_{n\to\infty} \sqrt n \left( \sqrt n-\sqrt{n-1} \right) =\infty \)}{divergentní a nemá nevlastní limitu}{}}
\LANGsk{
\Question{
Postupnosť \( \left( \sqrt n \left( \sqrt n-\sqrt{n-1} \right) \right)_{n=1}^{\infty} \) je:}
{konvergentná a platí: \( \lim\limits_{n\to\infty} \sqrt n \left( \sqrt n-\sqrt{n-1} \right) =\frac12 \)}{konvergentná a platí: \( \lim\limits_{n\to\infty} \sqrt n \left( \sqrt n-\sqrt{n-1} \right) =0 \)}{konvergentná a platí: \( \lim\limits_{n\to\infty} \sqrt n \left( \sqrt n-\sqrt{n-1} \right) =2 \)}{divergentná a platí: \( \lim\limits_{n\to\infty} \sqrt n \left( \sqrt n-\sqrt{n-1} \right) =\infty \)}{divergentná a nemá nevlastnú limitu}{}}
\LANGes{
\Question{
La sucesión \( \left( \sqrt n \left( \sqrt n-\sqrt{n-1} \right) \right)_{n=1}^{\infty} \) es:}
{convergente y \( \lim\limits_{n\to\infty} \sqrt n \left( \sqrt n-\sqrt{n-1} \right) =\frac12 \)}{convergente y \( \lim\limits_{n\to\infty} \sqrt n \left( \sqrt n-\sqrt{n-1} \right) =0 \)}{convergente y \( \lim\limits_{n\to\infty} \sqrt n \left( \sqrt n-\sqrt{n-1} \right) =2 \)}{divergente y \( \lim\limits_{n\to\infty} \sqrt n \left( \sqrt n-\sqrt{n-1} \right) =\infty \)}{divergente y no tiene un límite infinito.}{}}
\End

\Data\ID{1003047607}
\Updated{2020-07-15 03:07:12}
\Level{C}
\SubArea{Limit of a sequence}
\LANGen{
\Question{
Find the  limit  \( \lim\limits_{n\to\infty} n\left( \sqrt n-\sqrt{n-1} \right) \).}
{\( \infty \)}{\( \frac12 \)}{\( 0 \)}{\( 2 \)}{\( -\infty \)}{}}
\LANGpl{
\Question{
Wyznacz granicę dla  \( \lim\limits_{n\to\infty} n\left( \sqrt n-\sqrt{n-1} \right) \).}
{\( \infty \)}{\( \frac12 \)}{\( 0 \)}{\( 2 \)}{\( -\infty \)}{}}
\LANGcs{
\Question{
\( \lim\limits_{n\to\infty} n\left( \sqrt n-\sqrt{n-1} \right) \) je rovna:}
{\( \infty \)}{\( \frac12 \)}{\( 0 \)}{\( 2 \)}{\( -\infty \)}{}}
\LANGsk{
\Question{
\( \lim\limits_{n\to\infty} n\left( \sqrt n-\sqrt{n-1} \right) \) je rovná:}
{\( \infty \)}{\( \frac12 \)}{\( 0 \)}{\( 2 \)}{\( -\infty \)}{}}
\LANGes{
\Question{
Halla: \( \lim\limits_{n\to\infty} n\left( \sqrt n-\sqrt{n-1} \right) \).}
{\( \infty \)}{\( \frac12 \)}{\( 0 \)}{\( 2 \)}{\( -\infty \)}{}}
\End

\Data\ID{1003047608}
\Updated{2021-05-03 08:05:09}
\Level{C}
\SubArea{Limit of a sequence}
\LANGen{
\Question{
Choose the step to take first to efficiently evaluate the limit of the sequence \( \left( \frac{3n+2}{\sqrt{n^2-1}} \right)_{n=1}^{\infty} \).}
{We divide the numerator and the denominator by \( n \).}{We take \( \sqrt n \) outside brackets in the numerator and the denominator.}{We square the denominator.}{We divide the numerator by \( n \).}{We divide the denominator by \( n \).}{}}
\LANGpl{
\Question{
Wybierz pierwszy krok, który pozwoli skutecznie oszacować granicę ciągu \( \left( \frac{3n+2}{\sqrt{n^2-1}} \right)_{n=1}^{\infty} \).}
{Dzielimy licznik przez mianownik \( n \).}{Wyciągamy  \( \sqrt n \) poza nawias w liczniki i mianowniku.}{Podnosimy mianownik do kwadratu.}{Dzielimy licznik przez \( n \).}{Dzielimy mianownik przez \( n \).}{}}
\LANGcs{
\Question{
Vyberte vhodný postup pro výpočet limity posloupnosti \( \left( \frac{3n+2}{\sqrt{n^2-1}} \right)_{n=1}^{\infty} \).}
{Vydělíme čitatel i jmenovatel \( n \).}{Vytkneme v čitateli i jmenovateli \( \sqrt n \) .}{Umocníme jmenovatel.}{Vydělíme čitatel \( n \).}{Vydělíme jmenovatel \( n \).}{}}
\LANGsk{
\Question{
Vyberte vhodný postup pre výpočet limity postupnosti \( \left( \frac{3n+2}{\sqrt{n^2-1}} \right)_{n=1}^{\infty} \).}
{Vydelíme čitateľ i menovateľ \( n \).}{Vyjmeme v čitateli i menovateli \( \sqrt n \).}{Umocníme menovateľ.}{Vydelíme čitateľ \( n \).}{Vydelíme menovateľ \( n \).}{}}
\LANGes{
\Question{
Elige el primer paso más adecuado para calcular el límite de la siguiente sucesión:
 \( \left( \frac{3n+2}{\sqrt{n^2-1}} \right)_{n=1}^{\infty} \).}
{Dividir el numerador y el denominador por \( n \).}{Factorizar \( \sqrt n \) en el numerador y en el denominador.}{Elevar al cuadrado el denominador.}{Dividir el numerador por \( n \).}{Dividir el denominador por \( n \).}{}}
\End

\Data\ID{1003047609}
\Updated{2020-07-15 03:07:12}
\Level{C}
\SubArea{Limit of a sequence}
\LANGen{
\Question{
Find the  limit  \( \lim\limits_{n\to\infty}\frac{\sqrt{4n^2+3n+1}}{\sqrt{2n^2-5n-7}} \).}
{\( \sqrt2 \)}{\( 2 \)}{\( -\frac17 \)}{\( 0 \)}{\( \infty \)}{}}
\LANGpl{
\Question{
Wyznacz granicę dla  \( \lim\limits_{n\to\infty}\frac{\sqrt{4n^2+3n+1}}{\sqrt{2n^2-5n-7}} \).}
{\( \sqrt2 \)}{\( 2 \)}{\( -\frac17 \)}{\( 0 \)}{\( \infty \)}{}}
\LANGcs{
\Question{
\( \lim\limits_{n\to\infty}\frac{\sqrt{4n^2+3n+1}}{\sqrt{2n^2-5n-7}} \) je rovna:}
{\( \sqrt2 \)}{\( 2 \)}{\( -\frac17 \)}{\( 0 \)}{\( \infty \)}{}}
\LANGsk{
\Question{
\( \lim\limits_{n\to\infty}\frac{\sqrt{4n^2+3n+1}}{\sqrt{2n^2-5n-7}} \) je rovná:}
{\( \sqrt2 \)}{\( 2 \)}{\( -\frac17 \)}{\( 0 \)}{\( \infty \)}{}}
\LANGes{
\Question{
Halla:  \( \lim\limits_{n\to\infty}\frac{\sqrt{4n^2+3n+1}}{\sqrt{2n^2-5n-7}} \).}
{\( \sqrt2 \)}{\( 2 \)}{\( -\frac17 \)}{\( 0 \)}{\( \infty \)}{}}
\End

\Data\ID{1003047610}
\Updated{2020-07-15 03:07:12}
\Level{C}
\SubArea{Limit of a sequence}
\LANGen{
\Question{
Find the limit \( \lim\limits_{n\to\infty} \frac{\sqrt{4n+5}}{\sqrt{2n^3+3n^2-1}} \).}
{\( 0 \)}{\( \sqrt2 \)}{\( 2 \)}{\( \infty \)}{\( -5 \)}{}}
\LANGpl{
\Question{
Wyznacz granicę dla \( \lim\limits_{n\to\infty} \frac{\sqrt{4n+5}}{\sqrt{2n^3+3n^2-1}} \).}
{\( 0 \)}{\( \sqrt2 \)}{\( 2 \)}{\( \infty \)}{\( -5 \)}{}}
\LANGcs{
\Question{
\( \lim\limits_{n\to\infty} \frac{\sqrt{4n+5}}{\sqrt{2n^3+3n^2-1}} \) je rovna:}
{\( 0 \)}{\( \sqrt2 \)}{\( 2 \)}{\( \infty \)}{\( -5 \)}{}}
\LANGsk{
\Question{
\( \lim\limits_{n\to\infty} \frac{\sqrt{4n+5}}{\sqrt{2n^3+3n^2-1}} \) je rovná:}
{\( 0 \)}{\( \sqrt2 \)}{\( 2 \)}{\( \infty \)}{\( -5 \)}{}}
\LANGes{
\Question{
Halla \( \lim\limits_{n\to\infty} \frac{\sqrt{4n+5}}{\sqrt{2n^3+3n^2-1}} \).}
{\( 0 \)}{\( \sqrt2 \)}{\( 2 \)}{\( \infty \)}{\( -5 \)}{}}
\End

\Data\ID{1003047701}
\Updated{2020-07-15 03:07:12}
\Level{C}
\SubArea{Limit of a sequence}
\LANGen{
\Question{
Find the limit \(  \lim\limits_{n\to\infty}\frac{3+9+\dots+3^n}{4+16+\dots+4^n } \).}
{\( 0 \)}{\( \frac32 \)}{\( \infty \)}{\( 1 \)}{\( \frac34 \)}{}}
\LANGpl{
\Question{
Wyznacz granicę  \(  \lim\limits_{n\to\infty}\frac{3+9+\dots+3^n}{4+16+\dots+4^n } \).}
{\( 0 \)}{\( \frac32 \)}{\( \infty \)}{\( 1 \)}{\( \frac34 \)}{}}
\LANGcs{
\Question{
\(  \lim\limits_{n\to\infty}\frac{3+9+\dots+3^n}{4+16+\dots+4^n } \) je rovna:}
{\( 0 \)}{\( \frac32 \)}{\( \infty \)}{\( 1 \)}{\( \frac34 \)}{}}
\LANGsk{
\Question{
\(  \lim\limits_{n\to\infty}\frac{3+9+\dots+3^n}{4+16+\dots+4^n } \) je rovná:}
{\( 0 \)}{\( \frac32 \)}{\( \infty \)}{\( 1 \)}{\( \frac34 \)}{}}
\LANGes{
\Question{
Halla \(  \lim\limits_{n\to\infty}\frac{3+9+\dots+3^n}{4+16+\dots+4^n } \).}
{\( 0 \)}{\( \frac32 \)}{\( \infty \)}{\( 1 \)}{\( \frac34 \)}{}}
\End

\Data\ID{1003047702}
\Updated{2020-10-05 05:10:12}
\Level{C}
\SubArea{Limit of a sequence}
\LANGen{
\Question{
Find the limit:
\[ \lim\limits_{n\to\infty}\left(\frac13+\frac19+\dots+\frac1{3^n} \right) \].}
{\( \frac12 \)}{\( \frac13 \)}{\( \frac32 \)}{\( \infty \)}{\( \frac23 \)}{}}
\LANGpl{
\Question{
Wyznacz granicę \( \lim\limits_{n\to\infty}\left(\frac13+\frac19+\dots+\frac1{3^n} \right) \).}
{\( \frac12 \)}{\( \frac13 \)}{\( \frac32 \)}{\( \infty \)}{\( \frac23 \)}{}}
\LANGcs{
\Question{
Určete následující limitu.
\[ \lim\limits_{n\to\infty}\left(\frac13+\frac19+\dots+\frac1{3^n} \right) \]}
{\( \frac12 \)}{\( \frac13 \)}{\( \frac32 \)}{\( \infty \)}{\( \frac23 \)}{}}
\LANGsk{
\Question{
Určite následujúcu limitu.
\[ \lim\limits_{n\to\infty}\left(\frac13+\frac19+\dots+\frac1{3^n} \right) \]}
{\( \frac12 \)}{\( \frac13 \)}{\( \frac32 \)}{\( \infty \)}{\( \frac23 \)}{}}
\LANGes{
\Question{
Halla \( \lim\limits_{n\to\infty}\left(\frac13+\frac19+\dots+\frac1{3^n} \right) \).}
{\( \frac12 \)}{\( \frac13 \)}{\( \frac32 \)}{\( \infty \)}{\( \frac23 \)}{}}
\End

\Data\ID{1003047703}
\Updated{2020-07-15 03:07:12}
\Level{C}
\SubArea{Limit of a sequence}
\LANGen{
\Question{
Find the limit \( \lim\limits_{n\to\infty}\frac{\frac12+\frac14+\dots+\frac1{2^n}}{\frac13+\frac19+\dots+\frac1{3^n}} \).}
{\( 2 \)}{\( \frac23 \)}{\( \infty \)}{\( 0 \)}{\( \frac32 \)}{}}
\LANGpl{
\Question{
Wyznacz granicę  \( \lim\limits_{n\to\infty}\frac{\frac12+\frac14+\dots+\frac1{2^n}}{\frac13+\frac19+\dots+\frac1{3^n}} \).}
{\( 2 \)}{\( \frac23 \)}{\( \infty \)}{\( 0 \)}{\( \frac32 \)}{}}
\LANGcs{
\Question{
\( \lim\limits_{n\to\infty}\frac{\frac12+\frac14+\dots+\frac1{2^n}}{\frac13+\frac19+\dots+\frac1{3^n}} \) je rovna:}
{\( 2 \)}{\( \frac23 \)}{\( \infty \)}{\( 0 \)}{\( \frac32 \)}{}}
\LANGsk{
\Question{
\( \lim\limits_{n\to\infty}\frac{\frac12+\frac14+\dots+\frac1{2^n}}{\frac13+\frac19+\dots+\frac1{3^n}} \) je rovná:}
{\( 2 \)}{\( \frac23 \)}{\( \infty \)}{\( 0 \)}{\( \frac32 \)}{}}
\LANGes{
\Question{
Halla \( \lim\limits_{n\to\infty}\frac{\frac12+\frac14+\dots+\frac1{2^n}}{\frac13+\frac19+\dots+\frac1{3^n}} \).}
{\( 2 \)}{\( \frac23 \)}{\( \infty \)}{\( 0 \)}{\( \frac32 \)}{}}
\End

\Data\ID{1003047704}
\Updated{2020-07-15 04:07:51}
\Level{C}
\SubArea{Limit of a sequence}
\LANGen{
\Question{
Find the limit.  \[ \lim\limits_{n\to\infty}\frac{1^2+2^2+\dots+n^2}{n^2+7n-3} \]
Hint: \( 1^2+2^2+\cdots +n^2=\frac16 n(n+1)(2n+1) \).}
{\( \infty \)}{\( 0 \)}{\( \frac13 \)}{\( -\frac13 \)}{\( \frac16 \)}{}}
\LANGpl{
\Question{
Wyznacz granicę.  \[ \lim\limits_{n\to\infty}\frac{1^2+2^2+\dots+n^2}{n^2+7n-3} \]
Wskazówka: \( 1^2+2^2+\cdots +n^2=\frac16 n(n+1)(2n+1) \).}
{\( \infty \)}{\( 0 \)}{\( \frac13 \)}{\( -\frac13 \)}{\( \frac16 \)}{}}
\LANGcs{
\Question{
Vypočítejte limitu. \[ \lim\limits_{n\to\infty}\frac{1^2+2^2+\dots+n^2}{n^2+7n-3} \]
Nápověda: \( 1^2+2^2+\cdots +n^2=\frac16 n(n+1)(2n+1) \).}
{\( \infty \)}{\( 0 \)}{\( \frac13 \)}{\( -\frac13 \)}{\( \frac16 \)}{}}
\LANGsk{
\Question{
Vypočítajte limitu. \[ \lim\limits_{n\to\infty}\frac{1^2+2^2+\dots+n^2}{n^2+7n-3} \]
Nápoveda: \( 1^2+2^2+\cdots +n^2=\frac16 n(n+1)(2n+1) \).}
{\( \infty \)}{\( 0 \)}{\( \frac13 \)}{\( -\frac13 \)}{\( \frac16 \)}{}}
\LANGes{
\Question{
Halla:  \[ \lim\limits_{n\to\infty}\frac{1^2+2^2+\dots+n^2}{n^2+7n-3} \]
Sugerencia: \( 1^2+2^2+\cdots +n^2=\frac16 n(n+1)(2n+1) \).}
{\( \infty \)}{\( 0 \)}{\( \frac13 \)}{\( -\frac13 \)}{\( \frac16 \)}{}}
\End

\Data\ID{9000063610}
\Updated{2020-07-15 03:07:37}
\Level{C}
\SubArea{Limit of a sequence}
\LANGen{
\Question{
Find the following limit.
\[ 
 \lim _{n\to \infty } \frac{3 + 6 + 9 +\cdots +3n} 
{6 + 12 + 18 +\cdots +6n}
\]}
{\(\frac{1}
{2}\)}{\(0\)}{\(1\)}{\(\infty \)}{}{}}
\LANGpl{
\Question{
Wyznacz granicę ciągu.
\[ 
 \lim _{n\to \infty } \frac{3 + 6 + 9 +\cdots +3n} 
{6 + 12 + 18 +\cdots +6n}
\]}
{\(\frac{1}
{2}\)}{\(0\)}{\(1\)}{\(\infty \)}{}{}}
\LANGcs{
\Question{
\(\lim _{n\to \infty } \frac{3+6+9+\cdots +3n} 
{6+12+18+\cdots +6n}\) je
rovna:}
{\(\frac{1}
{2}\)}{\(0\)}{\(1\)}{\(\infty \)}{}{}}
\LANGsk{
\Question{
\(\lim _{n\to \infty } \frac{3+6+9+\cdots +3n} 
{6+12+18+\cdots +6n}\) je
rovná:}
{\(\frac{1}
{2}\)}{\(0\)}{\(1\)}{\(\infty \)}{}{}}
\LANGes{
\Question{
Halla:
\[ 
 \lim _{n\to \infty } \frac{3 + 6 + 9 +\cdots +3n} 
{6 + 12 + 18 +\cdots +6n}
\]}
{\(\frac{1}
{2}\)}{\(0\)}{\(1\)}{\(\infty \)}{}{}}
\End

\Data\ID{9000064001}
\Updated{2020-07-15 03:07:37}
\Level{C}
\SubArea{Limit of a sequence}
\LANGen{
\Question{
Consider the sequence 
\[ 
 \left (\frac{(-1)^{n}}
 {n} + 3\right )_{n=1}^{\infty }
\]
and its limit \(L\). How many terms
of the sequence differ from \(L\) 
by more than \(\frac{1}
{50}\)?}
{\(49\)}{\(50\)}{\(10\)}{\(100\)}{}{}}
\LANGpl{
\Question{
Rozważ ciąg
\[ 
 \left (\frac{(-1)^{n}}
 {n} + 3\right )_{n=1}^{\infty }
\]
i jego granicę \(L\). Ile wyrazów
różni się od  \(L\) 
o więcej niż \(\frac{1}
{50}\)?}
{\(49\)}{\(50\)}{\(10\)}{\(100\)}{}{}}
\LANGcs{
\Question{
Je dána konvergentní posloupnost \(\left (\frac{(-1)^{n}}
 {n} + 3\right )_{n=1}^{\infty }\).
Kolik členů této posloupnosti se liší od její limity o více než
\(\frac{1}
{50}\)?}
{\(49\)}{\(50\)}{\(10\)}{\(100\)}{}{}}
\LANGsk{
\Question{
Je daná konvergentná postupnosť \(\left (\frac{(-1)^{n}}
 {n} + 3\right )_{n=1}^{\infty }\).
Koľko členov tejto postupnosti sa líši od jej limity o viac než
\(\frac{1}
{50}\)?}
{\(49\)}{\(50\)}{\(10\)}{\(100\)}{}{}}
\LANGes{
\Question{
Dada la sucesión convergente: \(\left (\frac{(-1)^{n}}
 {n} + 3\right )_{n=1}^{\infty }\).
¿Cuántos términos de la sucesión difieren del límite en más de
\(\frac{1}
{50}\)?}
{\(49\)}{\(50\)}{\(10\)}{\(100\)}{}{}}
\End

\Data\ID{9000064002}
\Updated{2020-07-15 03:07:37}
\Level{C}
\SubArea{Limit of a sequence}
\LANGen{
\Question{
Consider the convergent sequence 
\[ 
 \left ( \frac{5 - n}
{2n - 1}\right )_{n=1}^{\infty }
\]
and its limit \(L\).
Find the index of the first term of the sequence which differs from
\(L\) by less
than \(\frac{1}
{100}\).}
{\(226\)}{\(450\)}{\(452\)}{\(451\)}{\(225\)}{}}
\LANGpl{
\Question{
Rozważ ciąg zbieżny
\[ 
 \left ( \frac{5 - n}
{2n - 1}\right )_{n=1}^{\infty }
\]
i jego granicę \(L\).
Wyznacz składnik pierwszego wyrazu ciągu, który różni się od 
\(L\) o mniej 
niż \(\frac{1}
{100}\).}
{\(226\)}{\(450\)}{\(452\)}{\(451\)}{\(225\)}{}}
\LANGcs{
\Question{
Je dána konvergentní posloupnost \(\left ( \frac{5-n}
{2n-1}\right )_{n=1}^{\infty }\).
Kterým členem počínaje se bude jeho hodnota lišit od limity o méně než
\(\frac{1}
{100}\)?}
{\(226\)}{\(450\)}{\(452\)}{\(451\)}{\(225\)}{}}
\LANGsk{
\Question{
Je daná konvergentná postupnosť \(\left ( \frac{5-n}
{2n-1}\right )_{n=1}^{\infty }\).
Ktorým členom počnúc sa bude jeho hodnota líšiť od limity o menej než
\(\frac{1}
{100}\)?}
{\(226\)}{\(450\)}{\(452\)}{\(451\)}{\(225\)}{}}
\LANGes{
\Question{
Dada la sucesión convergente:  \(\left ( \frac{5-n}
{2n-1}\right )_{n=1}^{\infty }\).
Halla el primer término de la sucesión que difiera del límite en menos de
\(\frac{1}
{100}\).}
{\(226\)}{\(450\)}{\(452\)}{\(451\)}{\(225\)}{}}
\End

\Data\ID{9000064003}
\Updated{2020-07-15 03:07:37}
\Level{C}
\SubArea{Limit of a sequence}
\LANGen{
\Question{
Consider the convergent sequence 
\[ 
 (a_{n})_{n=1}^{\infty } = \left (\frac{4n^{2} + 3n - 250}
 {2n^{2}} \right )_{n=1}^{\infty }
\]
and its limit \(L\). Find the
maximal difference between \(L\) 
and the subsequence \((a_{n})_{n=250}^{\infty }\).
(In other words, find the maximal difference between
\(L\) and the terms of the
sequence starting at \(a_{250}\).)}
{\(0.004\)}{\(0.04\)}{\(0.504\)}{\(0.54\)}{}{}}
\LANGpl{
\Question{
Rozważ ciąg zbieżny
\[ 
 (a_{n})_{n=1}^{\infty } = \left (\frac{4n^{2} + 3n - 250}
 {2n^{2}} \right )_{n=1}^{\infty }
\]
i jego granicę \(L\). Wyznacz
maksymalną różnicę pomiędzy \(L\) 
i podciąg  \((a_{n})_{n=250}^{\infty }\).
(Innymi słowy wyznacz maksymalną różnicę pomiędzy
\(L\) a składnikami ciągu
zaczynając od \(a_{250}\).)}
{\(0.004\)}{\(0.04\)}{\(0.504\)}{\(0.54\)}{}{}}
\LANGcs{
\Question{
Je dána konvergentní posloupnost \(\left (\frac{4n^{2}+3n-250}
 {2n^{2}} \right )_{n=1}^{\infty }\).
Určete maximální odchylku \(a_{n},n\geq 250\) 
od limity dané posloupnosti. (O kolik nejvíce se liší
\(a_{250}\) a
další členy posloupnosti od její limity?)}
{\(0{,}004\)}{\(0{,}04\)}{\(0{,}504\)}{\(0{,}54\)}{}{}}
\LANGsk{
\Question{
Je daná konvergentná postupnosť \(\left (\frac{4n^{2}+3n-250}
 {2n^{2}} \right )_{n=1}^{\infty }\).
Určte maximálnu odchýlku \(a_{n},n\geq 250\) 
od limity danej postupnosti. (O koľko najviac sa líši
\(a_{250}\) a
ďalší členy postupnosti od jej limity?)}
{\(0{,}004\)}{\(0{,}04\)}{\(0{,}504\)}{\(0{,}54\)}{}{}}
\LANGes{
\Question{
Dada la sucesión convergente
\[ 
 (a_{n})_{n=1}^{\infty } = \left (\frac{4n^{2} + 3n - 250}
 {2n^{2}} \right )_{n=1}^{\infty }
\]
y su límite \(L\). Halla la diferencia máxima entre \(L\) 
y la sucesión \((a_{n})_{n=250}^{\infty }\).
(es decir, halla la diferencia máxima entre
\(L\) y los términos de la sucesión que comienza en \(a_{250}\).)}
{\(0.004\)}{\(0.04\)}{\(0.504\)}{\(0.54\)}{}{}}
\End

\Data\ID{9000064008}
\Updated{2020-07-15 03:07:37}
\Level{C}
\SubArea{Limit of a sequence}
\LANGen{
\Question{
Find the limit of the following sequence. 
\[ 
 {\Bigl (\frac{(n^{2} + 2n + 1)^{n}}
 {n^{2n}} \Bigr )}_{n=1}^{\infty }
\]
Hint: The limit of the sequence \({\bigl ({\bigl (1 + \frac{1} 
{n}\bigr )}^{n}\bigr )}_{n=1}^{\infty }\) 
is the Euler number \(\mathrm{e}\).}
{\(\mathrm{e}^{2}\)}{\(2\mathrm{e}\)}{\(\mathrm{e} + 2\)}{\(\infty \)}{}{}}
\LANGpl{
\Question{
Wyznacz granicę ciągu.
\[ 
 {\Bigl (\frac{(n^{2} + 2n + 1)^{n}}
 {n^{2n}} \Bigr )}_{n=1}^{\infty }
\]
Wskazówka: Granicą ciągu \({\bigl ({\bigl (1 + \frac{1} 
{n}\bigr )}^{n}\bigr )}_{n=1}^{\infty }\) 
jest liczba Eulera \(\mathrm{e}\).}
{\(\mathrm{e}^{2}\)}{\(2\mathrm{e}\)}{\(\mathrm{e} + 2\)}{\(\infty \)}{}{}}
\LANGcs{
\Question{
Určete limitu posloupnosti. \[{\Bigl (\frac{(n^{2}+2n+1)^{n}}
 {n^{2n}} \Bigr )}_{n=1}^{\infty }\] Nápověda: Posloupnost \({\bigl ({\bigl (1 + \frac{1} 
{n}\bigr )}^{n}\bigr )}_{n=1}^{\infty }\) 
je konvergentní a její limita je Eulerovo číslo
\(\mathrm{e}\).}
{\(\mathrm{e}^{2}\)}{\(2\mathrm{e}\)}{\(\mathrm{e} + 2\)}{\(\infty\)}{}{}}
\LANGsk{
\Question{
Určte limitu postupnosti. \[{\Bigl (\frac{(n^{2}+2n+1)^{n}}
 {n^{2n}} \Bigr )}_{n=1}^{\infty }\] Nápoveda: Postupnosť \({\bigl ({\bigl (1 + \frac{1} 
{n}\bigr )}^{n}\bigr )}_{n=1}^{\infty }\) 
je konvergentné a jej limita je Eulerove číslo
\(\mathrm{e}\).}
{\(\mathrm{e}^{2}\)}{\(2\mathrm{e}\)}{\(\mathrm{e} + 2\)}{\(\infty\)}{}{}}
\LANGes{
\Question{
Halla:
\[ 
 {\Bigl (\frac{(n^{2} + 2n + 1)^{n}}
 {n^{2n}} \Bigr )}_{n=1}^{\infty }
\]
Sugerencia: El límite de la sucesión \({\bigl ({\bigl (1 + \frac{1} 
{n}\bigr )}^{n}\bigr )}_{n=1}^{\infty }\) 
es el número de Euler \(\mathrm{e}\).}
{\(\mathrm{e}^{2}\)}{\(2\mathrm{e}\)}{\(\mathrm{e} + 2\)}{\(\infty \)}{}{}}
\End

\Data\ID{9000064009}
\Updated{2020-07-15 03:07:37}
\Level{C}
\SubArea{Limit of a sequence}
\LANGen{
\Question{
Find the limit of the following sequence. 
\[ 
 {\Bigl ({\Bigl (\frac{\root{n}\of{2}}
 {n} + \root{n}\of{2}\Bigr )}^{n}\Bigr )}_{
n=1}^{\infty }
\]
Hint: The limit of the sequence \({\bigl ({\bigl (1 + \frac{1} 
{n}\bigr )}^{n}\bigr )}_{n=1}^{\infty }\) 
is the Euler number \(\mathrm{e}\).}
{\(2\mathrm{e}\)}{\(\mathrm{e}^{2}\)}{\(\mathrm{e} + 2\)}{\(\infty \)}{}{}}
\LANGpl{
\Question{
Wyznacz granicę ciągu.
\[ 
 {\Bigl ({\Bigl (\frac{\root{n}\of{2}}
 {n} + \root{n}\of{2}\Bigr )}^{n}\Bigr )}_{
n=1}^{\infty }
\]
Wskazówka: Granicą ciągu \({\bigl ({\bigl (1 + \frac{1} 
{n}\bigr )}^{n}\bigr )}_{n=1}^{\infty }\) 
jest liczba Eulera \(\mathrm{e}\).}
{\(2\mathrm{e}\)}{\(\mathrm{e}^{2}\)}{\(\mathrm{e} + 2\)}{\(\infty \)}{}{}}
\LANGcs{
\Question{
Určete limitu posloupnosti.
\[{\Bigl ({\Bigl (\frac{\root{n}\of{2}}
 {n} + \root{n}\of{2}\Bigr )}^{n}\Bigr )}_{
n=1}^{\infty }\] Nápověda: Posloupnost \({\bigl ({\bigl (1 + \frac{1} 
{n}\bigr )}^{n}\bigr )}_{n=1}^{\infty }\) 
je konvergentní a její limita je Eulerovo číslo
\(\mathrm{e}\).}
{\(2\mathrm{e}\)}{\(\mathrm{e}^{2}\)}{\(\mathrm{e} + 2\)}{\(\infty\)}{}{}}
\LANGsk{
\Question{
Určte limitu postupnosti.
\[{\Bigl ({\Bigl (\frac{\root{n}\of{2}}
 {n} + \root{n}\of{2}\Bigr )}^{n}\Bigr )}_{
n=1}^{\infty }\] Nápoveda: Postupnosť \({\bigl ({\bigl (1 + \frac{1} 
{n}\bigr )}^{n}\bigr )}_{n=1}^{\infty }\) 
je konvergentná a jej limita je Eulerove číslo
\(\mathrm{e}\).}
{\(2\mathrm{e}\)}{\(\mathrm{e}^{2}\)}{\(\mathrm{e} + 2\)}{\(\infty\)}{}{}}
\LANGes{
\Question{
Halla:
\[ 
 {\Bigl ({\Bigl (\frac{\root{n}\of{2}}
 {n} + \root{n}\of{2}\Bigr )}^{n}\Bigr )}_{
n=1}^{\infty }
\]
Sugerencia: El límite de la sucesión \({\bigl ({\bigl (1 + \frac{1} 
{n}\bigr )}^{n}\bigr )}_{n=1}^{\infty }\) 
es el número de Euler \(\mathrm{e}\).}
{\(2\mathrm{e}\)}{\(\mathrm{e}^{2}\)}{\(\mathrm{e} + 2\)}{\(\infty \)}{}{}}
\End

\Data\ID{9000064010}
\Updated{2020-07-15 03:07:37}
\Level{C}
\SubArea{Limit of a sequence}
\LANGen{
\Question{
Find the limit of the following sequence. 
\[ 
 {\Bigl ({\Bigl (\frac{2n + 1}
 {n} \Bigr )}^{n}\Bigr )}_{
n=1}^{\infty }
\]
Hint: The limit of the sequence \({\bigl ({\bigl (1 + \frac{1} 
{n}\bigr )}^{n}\bigr )}_{n=1}^{\infty }\) 
is the Euler number \(\mathrm{e}\).}
{\(\infty \)}{\(2\mathrm{e}\)}{\(\mathrm{e}^{2}\)}{\(\mathrm{e} + 2\)}{}{}}
\LANGpl{
\Question{
Wyznacz granicę następującego ciągu.
\[ 
 {\Bigl ({\Bigl (\frac{2n + 1}
 {n} \Bigr )}^{n}\Bigr )}_{
n=1}^{\infty }
\]
Wskazówka: Granicą ciągu \({\bigl ({\bigl (1 + \frac{1} 
{n}\bigr )}^{n}\bigr )}_{n=1}^{\infty }\) 
jest liczba Eulera \(\mathrm{e}\).}
{\(\infty \)}{\(2\mathrm{e}\)}{\(\mathrm{e}^{2}\)}{\(\mathrm{e} + 2\)}{}{}}
\LANGcs{
\Question{
Určete limitu posloupnosti. \[{\Bigl ({\Bigl (\frac{2n+1}
 {n} \Bigr )}^{n}\Bigr )}_{
n=1}^{\infty }\] Nápověda: Posloupnost \({\bigl ({\bigl (1 + \frac{1} 
{n}\bigr )}^{n}\bigr )}_{n=1}^{\infty }\) 
je konvergentní a její limita je Eulerovo číslo
\(\mathrm{e}\).}
{\(\infty\)}{\(2\mathrm{e}\)}{\(\mathrm{e}^{2}\)}{\(\mathrm{e} + 2\)}{}{}}
\LANGsk{
\Question{
Určte limitu postupnosti. \[{\Bigl ({\Bigl (\frac{2n+1}
 {n} \Bigr )}^{n}\Bigr )}_{
n=1}^{\infty }\] Nápoveda: Postupnosť \({\bigl ({\bigl (1 + \frac{1} 
{n}\bigr )}^{n}\bigr )}_{n=1}^{\infty }\) 
je konvergentná a jej limita je Eulerove číslo
\(\mathrm{e}\).}
{\(\infty\)}{\(2\mathrm{e}\)}{\(\mathrm{e}^{2}\)}{\(\mathrm{e} + 2\)}{}{}}
\LANGes{
\Question{
Halla:
\[ 
 {\Bigl ({\Bigl (\frac{2n + 1}
 {n} \Bigr )}^{n}\Bigr )}_{
n=1}^{\infty }
\]
Sugerencia: El límite de la sucesión \({\bigl ({\bigl (1 + \frac{1} 
{n}\bigr )}^{n}\bigr )}_{n=1}^{\infty }\) 
es el número de Euler \(\mathrm{e}\).}
{\(\infty \)}{\(2\mathrm{e}\)}{\(\mathrm{e}^{2}\)}{\(\mathrm{e} + 2\)}{}{}}
\End
